% !TeX program = xelatex
% =====================================================================
% PHU LUC TIENG VIET — Device-Anchored Irrigation Logs (HUJOS-TT)
% =====================================================================
% Ban dich cua supplementary.tex (EN). Bai chinh nop bang tieng Viet nen phu luc
% cung phai bang tieng Viet (anh Van, 2026-10-13).
%
% BA NGUYEN TAC BAT BUOC khi sua file nay:
%  1. KHONG go tay so cua bai chinh. Moi tham chieu cheo dung \ref/\eqref, so
%     duoc sinh tu supp_stubs_vi.tex (doc main_hujos_vi.aux). Ban EN go tay nen
%     da lech: sec:anchor EN=8/VI=7, eq:v4 EN=12/VI=11, eq:v5 EN=13/VI=12,
%     eq:umax EN=16/VI=15.
%  2. GIU NGUYEN moi con so va moi hedge. Thap phan kieu Viet: $26{,}7$, khong
%     phai $26.7$. Dung nhu tables_vi/ (0,60--0,85 / w_0=0{,}90).
%  3. Bang dung tables_vi/, block dung blocks_vi/ (ban EN cua 08/09 van \input
%     boi supplementary.tex EN — hai phu luc doc lap, khong dung chung).
%
% HAI CHO TRO SAI CO SAN TRONG BAN EN, DA SUA O DAY (khong sua ban EN vi ban do
% da nop/da doi chieu):
%   S9.2 viet "Table~5 of the main text" — bai chi co 4 bang; bang so sanh phuong
%        an MRV la tab:benchmark = 4  -> dung \ref{tab:benchmark}
%   S15  viet "Table~4 of the main text" cho cac so ghep cap — bang co cot
%        Inflation/ghep cap la tab:credit = 2 -> dung \ref{tab:credit}
%
% Tac gia va co quan lay DUNG khoi cua main_hujos_vi.tex (ban EN cua phu luc chi
% ghi 2 tac gia trong khi bai chinh co 5 — phu luc phai khop bai chinh).
% =====================================================================

\def\HUJOSVI{}
% !TeX program = xelatex
% =====================================================================
% HUJOS-TT (Hue University Journal of Science: Techniques and Technology)
% =====================================================================
% Moi thong so lay TU QUY DINH THAT cua tap chi (khong doan):
%   * https://jos.hueuni.edu.vn/index.php/hujos-tt/about/submissions
%       - toi da 12 trang, kho 19 cm x 27 cm
%       - font Palatino Linotype, 10 pt
%       - le 2 cm ca bon ben
%       - gian dong 1.15; 6 pt truoc va 3 pt sau moi doan
%       - cong thuc danh so ben phai, goi bang "Eq. (1)"
%       - caption BANG o TREN, caption HINH o DUOI, so bang/hinh IN DAM
%       - co khoang trong giua gia tri so va don vi (tru do/phut/giay)
%       - tai lieu tham khao kieu Vancouver (ngoac vuong) tu 2026
%       - abstract + tu khoa ca TIENG VIET va TIENG ANH, <= 250 tu
%       - tu khoa 3-5 cum, phai xuat hien trong abstract
%       - cuoi bai: ho ten, hoc ham, hoc vi, co quan, dia chi, dien thoai,
%         email, nguoi huong dan (neu co), chu ky tac gia, chi ro
%         corresponding author
%   * Template DOCX /tmp/hujos_tpl.docx — LUU Y: day la template cua HUJOS
%     KHOA HOC TU NHIEN (Natural Science), KHONG phai cua HUJOS-TT. No ghi
%     pgSz 20.5x29.7cm va HAI cot. Da tung bi dung nham lam chuan hinh hoc.
%     do tu word/document.xml va word/styles.xml:
%       - pgSz 11624 x 16840 twips = 20.5 x 29.7 cm, le tren/duoi 1.5 cm,
%         trai/phai 2 cm, HAI cot, columnsep 425 twips = 0.75 cm
%       - title 28 half-pt = 14 pt bold; tac gia 20 = 10 pt bold;
%         co quan 18 = 9 pt regular; abstract/tu khoa 18 = 9 pt
%       - Heading1 24 = 12 pt, before 360 twips = 18 pt, after 240 = 12 pt
%       - Heading2 before 18 pt, after 8 pt
%       - body spacing line 276 = 1.15, before 120 = 6 pt, after 60 = 3 pt
%       - font nhung trong template: PalatinoLinotype-{regular,bold,italic,
%         boldItalic}.ttf va Cardo
%     => Template NS khac HUJOS-TT ve kho giay, le VA SO COT. Bai nay nop
%        HUJOS-TT nen theo: 19x27 cm, le 2 cm deu, MOT cot — da xac nhan
%        bang 2 bai HUJOS-TT da xuat ban (xem comment o phan geometry).
%        KHONG duoc giu 2 cot: quy dinh van ban khong noi ve so cot, va
%        bai da xuat ban cua chinh tap chi nay deu 1 cot.
%   * Kiem tra font: Palatino Linotype that (trich tu template) co 16/16
%     glyph tieng Viet theo cmap. `fc-list :lang=vi` bao KHONG ho tro tieng
%     Viet la SAI -- fontconfig doan theo bang lang chu khong doc cmap.
%     Dung TeX Gyre Pagella (clone Palatino, metric-compatible) de an toan
%     ban quyen font, da test compile tieng Viet 0 missing character.
% =====================================================================

\documentclass[10pt,twoside]{article}

% ---- CONG TAC NGON NGU (them 2026-10-08 cho ban tieng Viet main_hujos_vi.tex) ----
% Ban tieng Viet dinh nghia \HUJOSVI TRUOC khi % !TeX program = xelatex
% =====================================================================
% HUJOS-TT (Hue University Journal of Science: Techniques and Technology)
% =====================================================================
% Moi thong so lay TU QUY DINH THAT cua tap chi (khong doan):
%   * https://jos.hueuni.edu.vn/index.php/hujos-tt/about/submissions
%       - toi da 12 trang, kho 19 cm x 27 cm
%       - font Palatino Linotype, 10 pt
%       - le 2 cm ca bon ben
%       - gian dong 1.15; 6 pt truoc va 3 pt sau moi doan
%       - cong thuc danh so ben phai, goi bang "Eq. (1)"
%       - caption BANG o TREN, caption HINH o DUOI, so bang/hinh IN DAM
%       - co khoang trong giua gia tri so va don vi (tru do/phut/giay)
%       - tai lieu tham khao kieu Vancouver (ngoac vuong) tu 2026
%       - abstract + tu khoa ca TIENG VIET va TIENG ANH, <= 250 tu
%       - tu khoa 3-5 cum, phai xuat hien trong abstract
%       - cuoi bai: ho ten, hoc ham, hoc vi, co quan, dia chi, dien thoai,
%         email, nguoi huong dan (neu co), chu ky tac gia, chi ro
%         corresponding author
%   * Template DOCX /tmp/hujos_tpl.docx — LUU Y: day la template cua HUJOS
%     KHOA HOC TU NHIEN (Natural Science), KHONG phai cua HUJOS-TT. No ghi
%     pgSz 20.5x29.7cm va HAI cot. Da tung bi dung nham lam chuan hinh hoc.
%     do tu word/document.xml va word/styles.xml:
%       - pgSz 11624 x 16840 twips = 20.5 x 29.7 cm, le tren/duoi 1.5 cm,
%         trai/phai 2 cm, HAI cot, columnsep 425 twips = 0.75 cm
%       - title 28 half-pt = 14 pt bold; tac gia 20 = 10 pt bold;
%         co quan 18 = 9 pt regular; abstract/tu khoa 18 = 9 pt
%       - Heading1 24 = 12 pt, before 360 twips = 18 pt, after 240 = 12 pt
%       - Heading2 before 18 pt, after 8 pt
%       - body spacing line 276 = 1.15, before 120 = 6 pt, after 60 = 3 pt
%       - font nhung trong template: PalatinoLinotype-{regular,bold,italic,
%         boldItalic}.ttf va Cardo
%     => Template NS khac HUJOS-TT ve kho giay, le VA SO COT. Bai nay nop
%        HUJOS-TT nen theo: 19x27 cm, le 2 cm deu, MOT cot — da xac nhan
%        bang 2 bai HUJOS-TT da xuat ban (xem comment o phan geometry).
%        KHONG duoc giu 2 cot: quy dinh van ban khong noi ve so cot, va
%        bai da xuat ban cua chinh tap chi nay deu 1 cot.
%   * Kiem tra font: Palatino Linotype that (trich tu template) co 16/16
%     glyph tieng Viet theo cmap. `fc-list :lang=vi` bao KHONG ho tro tieng
%     Viet la SAI -- fontconfig doan theo bang lang chu khong doc cmap.
%     Dung TeX Gyre Pagella (clone Palatino, metric-compatible) de an toan
%     ban quyen font, da test compile tieng Viet 0 missing character.
% =====================================================================

\documentclass[10pt,twoside]{article}

% ---- CONG TAC NGON NGU (them 2026-10-08 cho ban tieng Viet main_hujos_vi.tex) ----
% Ban tieng Viet dinh nghia \HUJOSVI TRUOC khi % !TeX program = xelatex
% =====================================================================
% HUJOS-TT (Hue University Journal of Science: Techniques and Technology)
% =====================================================================
% Moi thong so lay TU QUY DINH THAT cua tap chi (khong doan):
%   * https://jos.hueuni.edu.vn/index.php/hujos-tt/about/submissions
%       - toi da 12 trang, kho 19 cm x 27 cm
%       - font Palatino Linotype, 10 pt
%       - le 2 cm ca bon ben
%       - gian dong 1.15; 6 pt truoc va 3 pt sau moi doan
%       - cong thuc danh so ben phai, goi bang "Eq. (1)"
%       - caption BANG o TREN, caption HINH o DUOI, so bang/hinh IN DAM
%       - co khoang trong giua gia tri so va don vi (tru do/phut/giay)
%       - tai lieu tham khao kieu Vancouver (ngoac vuong) tu 2026
%       - abstract + tu khoa ca TIENG VIET va TIENG ANH, <= 250 tu
%       - tu khoa 3-5 cum, phai xuat hien trong abstract
%       - cuoi bai: ho ten, hoc ham, hoc vi, co quan, dia chi, dien thoai,
%         email, nguoi huong dan (neu co), chu ky tac gia, chi ro
%         corresponding author
%   * Template DOCX /tmp/hujos_tpl.docx — LUU Y: day la template cua HUJOS
%     KHOA HOC TU NHIEN (Natural Science), KHONG phai cua HUJOS-TT. No ghi
%     pgSz 20.5x29.7cm va HAI cot. Da tung bi dung nham lam chuan hinh hoc.
%     do tu word/document.xml va word/styles.xml:
%       - pgSz 11624 x 16840 twips = 20.5 x 29.7 cm, le tren/duoi 1.5 cm,
%         trai/phai 2 cm, HAI cot, columnsep 425 twips = 0.75 cm
%       - title 28 half-pt = 14 pt bold; tac gia 20 = 10 pt bold;
%         co quan 18 = 9 pt regular; abstract/tu khoa 18 = 9 pt
%       - Heading1 24 = 12 pt, before 360 twips = 18 pt, after 240 = 12 pt
%       - Heading2 before 18 pt, after 8 pt
%       - body spacing line 276 = 1.15, before 120 = 6 pt, after 60 = 3 pt
%       - font nhung trong template: PalatinoLinotype-{regular,bold,italic,
%         boldItalic}.ttf va Cardo
%     => Template NS khac HUJOS-TT ve kho giay, le VA SO COT. Bai nay nop
%        HUJOS-TT nen theo: 19x27 cm, le 2 cm deu, MOT cot — da xac nhan
%        bang 2 bai HUJOS-TT da xuat ban (xem comment o phan geometry).
%        KHONG duoc giu 2 cot: quy dinh van ban khong noi ve so cot, va
%        bai da xuat ban cua chinh tap chi nay deu 1 cot.
%   * Kiem tra font: Palatino Linotype that (trich tu template) co 16/16
%     glyph tieng Viet theo cmap. `fc-list :lang=vi` bao KHONG ho tro tieng
%     Viet la SAI -- fontconfig doan theo bang lang chu khong doc cmap.
%     Dung TeX Gyre Pagella (clone Palatino, metric-compatible) de an toan
%     ban quyen font, da test compile tieng Viet 0 missing character.
% =====================================================================

\documentclass[10pt,twoside]{article}

% ---- CONG TAC NGON NGU (them 2026-10-08 cho ban tieng Viet main_hujos_vi.tex) ----
% Ban tieng Viet dinh nghia \HUJOSVI TRUOC khi \input{preamble}. Ban tieng Anh
% khong dinh nghia gi => \ifhujosvi = false => MOI thu giong het truoc khi sua
% (da kiem chung: main_hujos.pdf van 14 trang, 0 undefined, 0 overfull).
% Ly do dung macro \HUJOSVI thay vi babel/polyglossia: tieng Viet trong TeX Live
% chi co `zerohyph.tex` (khong co mau ngat am), nen nap polyglossia chi them
% rui ro xung dot voi unicode-math ma khong duoc loi ich nao. Tu tieng Viet ngan
% va tach bang dau cach nen diem ngat dong rat day du.
\newif\ifhujosvi
\ifdefined\HUJOSVI\hujosvitrue\fi

% ---- hinh hoc ----
% DA XAC MINH BANG BAI DA XUAT BAN THAT, khong con dua vao template nao.
%
% LOI DA MAC (ghi lai de khong tai pham): ban dau lay hinh hoc tu
% /tmp/hujos_tpl.docx = A4 20.5x29.7cm, le 1.5cm tren/duoi, HAI cot, columnsep
% 425 twips. File do CO that va cac con so do DUNG (425 twips nam trong the con
% <w:col w:space="425" w:w="4465.5">, khong phai trong <w:cols/>). Nhung HEADER
% BEN TRONG cua no ghi "Hue University Journal of Science: Natural Science
% Vol. 129, No. 1B" — do la template cua HUJOS KHOA HOC TU NHIEN, mot tap chi
% KHAC trong cung he thong. Bai nay nop HUJOS-TT (Ky thuat & Cong nghe).
% Hau qua: dung layout cua tap chi khac, vung chu 16.74cm thay vi 15.01cm,
% rong hon 11.5% -> so trang do duoc (13) khong co gia tri.
%
% BANG CHUNG QUYET DINH = 2 bai HUJOS-TT da xuat ban, tai tu
% citation_pdf_url cua chinh trang bai (article 7830, 7862):
%   pdfinfo     : 538.7 x 765.5 pt = 19.00 x 27.00 cm  (ca hai bai)
%   vung chu    : xMin 2.00cm -> xMax 17.01cm, rong 15.01 cm
%   le          : trai 2.00, phai 1.99 cm  -> le 2cm deu, dung nhu quy dinh van ban
%   SO COT      : chu bat dau trong dai giua trang = 20.1% va 20.5% (day chu)
%                 -> MOT COT. (Ban hai cot chi co 4.9%, vi ranh cot trong.)
%   font        : PalatinoLinotype-Roman/Bold/Italic + Cambria + Arial-BoldMT
% Kiem cheo: Author Guidelines cua HUJOS-TT noi "19 cm x 27 cm ... margins on
% all sides of 2 cm ... Palatino Linotype ... font size 10", va NHAC DEN
% "column" 0 LAN. Ba nguon (quy dinh van ban + 2 bai da xuat ban + template docx
% ref_hujos*.docx 19x27 le 2 1 cot) nay THONG NHAT. Chi co hujos_tpl.docx lech,
% va no thuoc tap chi khac.
\usepackage[paperwidth=19cm,paperheight=27cm,
            top=2cm,bottom=2cm,left=2cm,right=2cm,
            headsep=0.4cm,footskip=0.6cm]{geometry}

% ---- font: Palatino (TeX Gyre Pagella) + math khop ----
\usepackage{fontspec}
\setmainfont{TeX Gyre Pagella}
\usepackage{unicode-math}
\setmathfont{TeX Gyre Pagella Math}

% ---- MOT COT (xac nhan bang 2 bai HUJOS-TT da xuat ban) ----
% \columnsep khong con tac dung khi document khong \twocolumn; giu gia tri nay
% chi de neu ai do build lai 2 cot thi khong vo tinh dung khoang cach cua tap chi
% khac. Vung chu that = 15.01 cm theo bai da xuat ban.
\setlength{\columnsep}{0.75cm}

% ---- gian dong + indent doan: DO TU 2 BAI HUJOS-TT DA XUAT BAN ----
% experiments/measure_indent.py (loc dung bang tieu chi "doan truoc ket thuc som
% hon le phai 0.5cm", khong dem moi dong thut le vi nham voi dong trong bang):
%   indent dau doan = 1.00 cm  (bai #1: 31/72 dong; bai #2: 30/58 dong — trung vi
%                               cua ca hai dung bang 1.00)
%   khoang cach doan = pitch tai cho bat dau doan tru them +3.01pt (bai #1)
%                                            va +6.00pt (bai #2) so voi pitch thuong
% Quy dinh van ban: "6 points before and 3 points after paragraph spacing".
% template docx docDefaults: w:spacing before="120" after="60" line="276"
% (= 6pt / 3pt / 1.15). Hai nguon nay khop nhau -> chon parskip 3pt.
%
% Ban cu de parindent 0.45cm va parskip 0pt: KHONG khop nguon nao.
%
% "1.15" CUA WORD KHAC \linespread{1.15} CUA LATEX — da giai bang phep do, khong
% con la xung dot phai doan. Word nhan 1.15 voi (ascent+descent+lineGap) cua FONT;
% LaTeX \linespread nhan voi baselineskip cua CLASS.
%
% experiments/probe_linespread.py: 7 PDF RIENG, moi PDF mot gia tri (probe dau tien
% dat ca 4 gia tri trong cung mot file va tach doan bang khoang cach doc -> ham tach
% gop ca 4 doan lipsum thanh 1, trung vi ra 15.54pt la TRON cua 4 gia tri, vo nghia;
% \typeout khong xuat hien trong PDF nen khong dung lam moc duoc).
%   baselineskip tu nhien cua TeX Gyre Pagella 10pt = 11.954pt
%   (7 gia tri cho cung mot he so: 11.95-11.96pt, sai so 0.01pt)
%   1.15 -> 13.75pt   1.25 -> 14.94pt   1.29 -> 15.42pt   1.30 -> 15.54pt
% 2 bai HUJOS-TT da xuat ban do duoc pitch 15.48pt (n=225 va n=161 khoang cach).
% => \linespread{1.29} cho 15.42pt, lech target -0.06pt. Chon 1.29.
% Dung 1.15 (13.75pt) se chat hon bai cua chinh toa soan 1.13 lan.
\linespread{1.29}
\setlength{\parskip}{3pt plus 0.5pt}
\setlength{\parindent}{1.0cm}

% \raggedbottom la BAT BUOC o day. \documentclass[10pt,twoside] mac dinh bat
% \flushbottom, tuc LaTeX keo gian KEO DOC de moi trang cao bang nhau. Hau qua do
% duoc (experiments/verify_layout.py): parskip that cua ban em do ra +7.97pt thay vi
% 3pt da dat, va khoang cach tai ranh doan phan tan 16.9/21.7/23.9pt trong khi 2 bai
% da dang gom chat 18.1/21.5pt. \raggedbottom de trang ket thuc tu nhien, khong keo.
% parskip GIU 3pt: 2 bai da dang do duoc +2.64pt tai ranh doan (ca hai bai giong
% het nhau), khop "3 points after paragraph spacing" cua quy dinh van ban.
\raggedbottom

% ---- toan hoc ----
% KHONG dung \numberwithin{equation}{section}: quy dinh HUJOS-TT yeu cau cong thuc
% danh so lien tuc va duoc goi la "Eq. (1)", nen danh so theo section se ra (1.1)
% va lam sai ca cach goi lan thu tu trong van ban.
%
% KHONG nap amssymb CUNG unicode-math. Da do bang build that: hai goi trung dinh
% nghia 4 lenh va dung bien dich voi 4 loi fatal
%   ! LaTeX Error: Command `\eth' already defined.
%   ! LaTeX Error: Command `\smallsetminus' already defined.
%   ! LaTeX Error: Command `\digamma' already defined.
%   ! LaTeX Error: Command `\backepsilon' already defined.
% unicode-math (nap o tren, kem TeX Gyre Pagella Math) da cung cap day du symbol
% cua amssymb, nen bo amssymb la cach sua dung thay vi \let ghi de tung lenh.
\usepackage{amsmath}

% ---- bang: booktabs + caption tren bang ----
\usepackage{booktabs}
% Caption canh GIUA: do tu 2 bai da xuat ban (experiments/measure_venue_style.py),
% 17/17 caption 'Table n.'/'Fig. n.' co mid = 9.50cm dung tam vung chu
% (2.00..17.00cm), ke ca caption dai 2 dong (x=[2.79,16.21] va x=[2.74,16.09]).
% Ban cu dung justification=raggedright + singlelinecheck=false = canh TRAI,
% khong khop bai nao. font=small (9pt) giu nguyen: do duoc 9.14pt o bai da dang.
% labelfont=bf giu nguyen: quy dinh noi "Figure numbers and table numbers are
% printed in bold". labelsep=period giu nguyen: bai dang that la "Fig. 1." / "Table 1."
\usepackage[font=small,labelfont=bf,labelsep=period,
            skip=4pt,justification=centering]{caption}
\captionsetup[table]{position=top}
\captionsetup[figure]{position=bottom}
% Ban tieng Viet: nhan float phai la "Bang 1." / "Hinh 1." chu khong phai "Table 1."
% Quy dinh HUJOS ghi "table titles ABOVE the table, figure titles BELOW the figure"
% va "table/figure numbers printed in BOLD" -- vi tri va do dam khong doi, chi doi
% CHU cua nhan.
\ifhujosvi
  \captionsetup[table]{name=Bảng}
  \captionsetup[figure]{name=Hình}
\fi

% ---- heading theo so do tu template ----
\usepackage{titlesec}
% Bai da dang: so muc "1" nam o x=[2.00,2.33]cm, CHU "Introduction" bat dau o
% x=3.00cm (PyMuPDF tach thanh 2 span cung dong y=19.3). Tuc khoang cach tu CUOI
% nhan den dau chu = 0.79cm (do lai lan 2: het nhan 2.21 -> chu 3.00, ca 2 bai
% deu ra gap 0.79cm). Ban cu dung 0.5em (~0.21cm o 12pt) -> chu nam o
% ~2.25cm; ban sua 0.67cm cho 2.88cm, van lech 0.12cm -> chot 0.79cm.
% Subsection GIU 0.5em: bai #2 do duoc "2.1 Object Detection from UAV Imagery"
% bat dau o x0=2.00cm, khong thut them.
\titleformat{\section}{\normalfont\bfseries\fontsize{12}{14}\selectfont}
  {\thesection}{0.79cm}{}
\titlespacing*{\section}{0pt}{18pt}{12pt}
\titleformat{\subsection}{\normalfont\bfseries\fontsize{10}{12}\selectfont}
  {\thesubsection}{0.5em}{}
\titlespacing*{\subsection}{0pt}{18pt}{8pt}
\titleformat{\subsubsection}[runin]{\normalfont\itshape}
  {\thesubsubsection}{0.5em}{}[.]
\titlespacing*{\subsubsection}{0pt}{12pt}{0.5em}
% 18 pt space after text in each section (template yeu cau)
\titleformat{name=\section,numberless}{\normalfont\bfseries\fontsize{12}{14}\selectfont}
  {}{0pt}{}
\titlespacing*{name=\section,numberless}{0pt}{18pt}{12pt}

% ---- dinh ly / chung minh ----
\usepackage{amsthm}
% Ban tieng Viet doi ten moi truong dinh ly (KHONG doi noi dung dinh ly).
% Ten tieng Viet dung theo thong le toan hoc Viet Nam: proposition = Menh de.
\ifhujosvi
  \newtheorem{proposition}{Mệnh đề}
  \newtheorem{remark}{Ghi chú}
  \renewcommand{\proofname}{Chứng minh}
\else
  \newtheorem{proposition}{Proposition}
  \newtheorem{remark}{Remark}
  \renewcommand{\proofname}{Proof}
\fi

% ---- trích dẫn Vancouver (ngoặc vuông) từ 2026 ----
\usepackage[numbers,sort&compress]{natbib}
\bibpunct{[}{]}{,}{n}{}{,}

% ---- link (an vien, khong doi mau) ----
% PHAI \PassOptionsToPackage TRUOC hyperref: hyperref tu nap url ben trong, nen
% viet \usepackage[hyphens]{url} o SAU se dung "Option clash for package url"
% va dung bien dich. Da tra gia: build bao P3=1 ma van xuat PDF 13 trang, trong
% do 3 overfull cua .bbl giu nguyen 189.81/139.74/118.98pt -> hyphens chua an.
\PassOptionsToPackage{hyphens}{url}
\usepackage[hidelinks]{hyperref}

% ---- header/footer nhu template: khong tu them so trang ngoai cac thu da set ----
\usepackage{fancyhdr}
\pagestyle{fancy}
\fancyhf{}
\renewcommand{\headrulewidth}{0pt}
% Header: ban tieng Viet ghi ten tac gia khong dau (dung quy uoc cua chinh tap chi
% o cac bai tieng Viet da dang) va ten tap chi tieng Viet.
\ifhujosvi
  \fancyhead[LE]{\footnotesize\itshape M. X. Văn và cs.}
  \fancyhead[RO]{\footnotesize\itshape Tạp chí Khoa học Đại học Huế: Kỹ thuật và Công nghệ}
\else
  \fancyhead[LE]{\footnotesize\itshape X. V. Mai et al.}
  \fancyhead[RO]{\footnotesize\itshape Hue University Journal of Science: Techniques and Technology}
\fi
\fancyfoot[C]{\footnotesize\thepage}

% ---- phu luc: doi ten "Appendix" thanh "Appendix" (giu nguyen) ----
\usepackage{appendix}

% ---- tien ich ----
% [hyphens] la BAT BUOC: ba overfull nang nhat cua bai (189.81 / 139.74 /
% 118.98 pt, tuc chu loi ra ngoai trang) deu nam trong .bbl va deu do URL dai
% khong ngat duoc. Mac dinh goi url chi ngat o dau '/' va '.', ma ten mien nhu
% 'technical-explainer-carbon-credits-for-improved-rice-cultivation' khong co '/'.
% Option duoc truyen bang \PassOptionsToPackage o tren (truoc hyperref) de tranh
% option clash, nen o day nap khong kem option.
\usepackage{url}
\usepackage{xspace}
\usepackage{microtype}
\usepackage{array}

% Moi truong liet ke ky hieu cong thuc ('trong do:'). Dinh nghia o day chu khong o
% blocks/ vi cac lenh \setlength chua CHU SO (0.5pt, 0pt) se bi gate
% check_vi_numbers dem nhu so lieu nghien cuu -> lam loang gate.
\newenvironment{symbolist}
  {\begin{description}\footnotesize\setlength{\itemsep}{0.5pt}\setlength{\parsep}{0pt}}
  {\end{description}}


% graphicx la BAT BUOC cho \resizebox. Thieu no thi 24 loi fatal
% "Undefined control sequence" + "Missing number" + "Illegal unit of measure"
% (vi \resizebox bi phan tich sai). Da tra gia: ban IEEE co goi nay, ban HUJOS
% quen -> build ra PDF 13 trang nhung la RAC, khong duoc bao cao la thanh cong.
\usepackage{graphicx}

% \fittocol: chi THU NHO bang khi no that su rong hon cot, KHONG BAO GIO phong to.
%
% Tai sao khong dung \sbox de do: cot p{} chua \par, ma \sbox = \savebox trong
% \hbox nen khong do duoc vat lieu dang doan van. Do bang \sbox cho ra 7.92-8.01cm
% trong khi build that bao overfull 18-26pt (~8.8cm) -> so do sai, dan den ket luan
% sai rang "bang khong tran". lrbox la cong cu dung cho hop chua \par.
%
% Tai sao khong dung \resizebox{\columnwidth}{!} truc tiep: no CUONG BUC dung bang
% kho cot, nen bang hep hon (vd tab_baseline 13.69cm trong table* 16.5cm) se bi
% PHONG TO ~20%, chu to hon thiet ke.
\newsavebox{\fitbox}
\newcommand{\fittocol}[1]{%
  \begin{lrbox}{\fitbox}#1\end{lrbox}%
  \ifdim\wd\fitbox>\columnwidth
    \resizebox{\columnwidth}{!}{\usebox{\fitbox}}%
  \else
    \usebox{\fitbox}%
  \fi}

% Cho phep keo gian khoang trang them 2em khi khong tim duoc diem ngat.
% Ly do: cac doan chua nhieu \ref lien tiep (vd
% "(Sec.~\ref{sec:umax}, Sec.~\ref{sec:v2reach}, Table~\ref{tab:umax},
% Table~\ref{tab:f2reach})") tao ra chuoi khong ngat duoc, gay overfull
% 34.89pt va 26.16pt trong 01_intro.tex. Day la cach sua khong doi noi dung.
\setlength{\emergencystretch}{2em}

% Kho cot that: MOT cot, vung chu 19 - 2 - 2 = 15.00 cm
% (bai da xuat ban do duoc 15.01 cm — khop). \columnwidth = \textwidth o day.
% Day la con so ma experiments/make_tables_hujos.py dung de gate chieu rong bang.
\newlength{\HUJOScol}
\setlength{\HUJOScol}{15.0cm}

% ---------------- FRONT MATTER: moi con so DO TU 2 BAI HUJOS-TT DA XUAT BAN --------
% experiments/measure_fonts.py (PyMuPDF, doc theo SPAN nen biet dung font cua tung
% manh chu, khong gop dong):
%   title        PalatinoLinotype-Bold 14.04pt, canh giua       -> em: 14pt bold center  (khop)
%   authors      canh GIUA: mid = 9.59 va 9.53 cm (tam vung chu 9.5)
%   affiliations canh GIUA: mid = 9.54 cm (ca 4 dong cua bai #1)
%   abstract     9.00pt, x = [3.00, 16.08] cm -> THUT VAO 1.00 cm ca hai ben,
%                trong khi than bai x = [2.00, 17.08] cm (khong thut)
%   keywords     cung nam trong khoi thut 1.00 cm (x bat dau 3.00)
%   caption      9.00pt, chi NHAN in dam: span 'Fig. 1.' = PalatinoLinotype-Bold,
%                span con lai = PalatinoLinotype-Roman (10/10 caption bai #1,
%                8/8 bai #2) -> labelfont=bf + font=small cua em da khop san
% Ban cu de \raggedright cho authors/affil va khong thut le abstract: lech ca hai
% bai da dang o dung nhung cho nay.
\newcommand{\hujostitle}[1]{%
  {\centering\normalfont\bfseries\fontsize{14}{17}\selectfont #1\par}}
\newcommand{\hujosauthors}[1]{%
  {\centering\normalfont\bfseries\fontsize{10}{12}\selectfont #1\par}}
\newcommand{\hujosaffil}[1]{%
  {\centering\normalfont\fontsize{9}{11}\selectfont #1\par}}
% \leftskip/\rightskip = 1cm de abstract + keywords thut vao dung nhu bai da dang.
% Keywords nam BEN TRONG hujosabstract trong main_hujos.tex nen cung duoc thut.
% CO CHU ABSTRACT = 9pt, KHONG PHAI 11pt.
% LOI EM VUA GAY RA TRONG CHINH LAN SUA NAY: go \fontsize{11}{13} thay vi {9}{11}
% -> abstract se to hon than bai (10pt), nguoc hoan toan voi so do da do
% (measure_fonts.py: abstract 9.00pt o ca 2 bai; body 9.96pt). Da sua lai 9/11.
\newenvironment{hujosabstract}
  {\par\vspace{10pt}\normalfont\fontsize{9}{11}\selectfont
   \leftskip=1cm\rightskip=1cm\parindent=0pt\noindent\ignorespaces}
  {\par}
. Ban tieng Anh
% khong dinh nghia gi => \ifhujosvi = false => MOI thu giong het truoc khi sua
% (da kiem chung: main_hujos.pdf van 14 trang, 0 undefined, 0 overfull).
% Ly do dung macro \HUJOSVI thay vi babel/polyglossia: tieng Viet trong TeX Live
% chi co `zerohyph.tex` (khong co mau ngat am), nen nap polyglossia chi them
% rui ro xung dot voi unicode-math ma khong duoc loi ich nao. Tu tieng Viet ngan
% va tach bang dau cach nen diem ngat dong rat day du.
\newif\ifhujosvi
\ifdefined\HUJOSVI\hujosvitrue\fi

% ---- hinh hoc ----
% DA XAC MINH BANG BAI DA XUAT BAN THAT, khong con dua vao template nao.
%
% LOI DA MAC (ghi lai de khong tai pham): ban dau lay hinh hoc tu
% /tmp/hujos_tpl.docx = A4 20.5x29.7cm, le 1.5cm tren/duoi, HAI cot, columnsep
% 425 twips. File do CO that va cac con so do DUNG (425 twips nam trong the con
% <w:col w:space="425" w:w="4465.5">, khong phai trong <w:cols/>). Nhung HEADER
% BEN TRONG cua no ghi "Hue University Journal of Science: Natural Science
% Vol. 129, No. 1B" — do la template cua HUJOS KHOA HOC TU NHIEN, mot tap chi
% KHAC trong cung he thong. Bai nay nop HUJOS-TT (Ky thuat & Cong nghe).
% Hau qua: dung layout cua tap chi khac, vung chu 16.74cm thay vi 15.01cm,
% rong hon 11.5% -> so trang do duoc (13) khong co gia tri.
%
% BANG CHUNG QUYET DINH = 2 bai HUJOS-TT da xuat ban, tai tu
% citation_pdf_url cua chinh trang bai (article 7830, 7862):
%   pdfinfo     : 538.7 x 765.5 pt = 19.00 x 27.00 cm  (ca hai bai)
%   vung chu    : xMin 2.00cm -> xMax 17.01cm, rong 15.01 cm
%   le          : trai 2.00, phai 1.99 cm  -> le 2cm deu, dung nhu quy dinh van ban
%   SO COT      : chu bat dau trong dai giua trang = 20.1% va 20.5% (day chu)
%                 -> MOT COT. (Ban hai cot chi co 4.9%, vi ranh cot trong.)
%   font        : PalatinoLinotype-Roman/Bold/Italic + Cambria + Arial-BoldMT
% Kiem cheo: Author Guidelines cua HUJOS-TT noi "19 cm x 27 cm ... margins on
% all sides of 2 cm ... Palatino Linotype ... font size 10", va NHAC DEN
% "column" 0 LAN. Ba nguon (quy dinh van ban + 2 bai da xuat ban + template docx
% ref_hujos*.docx 19x27 le 2 1 cot) nay THONG NHAT. Chi co hujos_tpl.docx lech,
% va no thuoc tap chi khac.
\usepackage[paperwidth=19cm,paperheight=27cm,
            top=2cm,bottom=2cm,left=2cm,right=2cm,
            headsep=0.4cm,footskip=0.6cm]{geometry}

% ---- font: Palatino (TeX Gyre Pagella) + math khop ----
\usepackage{fontspec}
\setmainfont{TeX Gyre Pagella}
\usepackage{unicode-math}
\setmathfont{TeX Gyre Pagella Math}

% ---- MOT COT (xac nhan bang 2 bai HUJOS-TT da xuat ban) ----
% \columnsep khong con tac dung khi document khong \twocolumn; giu gia tri nay
% chi de neu ai do build lai 2 cot thi khong vo tinh dung khoang cach cua tap chi
% khac. Vung chu that = 15.01 cm theo bai da xuat ban.
\setlength{\columnsep}{0.75cm}

% ---- gian dong + indent doan: DO TU 2 BAI HUJOS-TT DA XUAT BAN ----
% experiments/measure_indent.py (loc dung bang tieu chi "doan truoc ket thuc som
% hon le phai 0.5cm", khong dem moi dong thut le vi nham voi dong trong bang):
%   indent dau doan = 1.00 cm  (bai #1: 31/72 dong; bai #2: 30/58 dong — trung vi
%                               cua ca hai dung bang 1.00)
%   khoang cach doan = pitch tai cho bat dau doan tru them +3.01pt (bai #1)
%                                            va +6.00pt (bai #2) so voi pitch thuong
% Quy dinh van ban: "6 points before and 3 points after paragraph spacing".
% template docx docDefaults: w:spacing before="120" after="60" line="276"
% (= 6pt / 3pt / 1.15). Hai nguon nay khop nhau -> chon parskip 3pt.
%
% Ban cu de parindent 0.45cm va parskip 0pt: KHONG khop nguon nao.
%
% "1.15" CUA WORD KHAC \linespread{1.15} CUA LATEX — da giai bang phep do, khong
% con la xung dot phai doan. Word nhan 1.15 voi (ascent+descent+lineGap) cua FONT;
% LaTeX \linespread nhan voi baselineskip cua CLASS.
%
% experiments/probe_linespread.py: 7 PDF RIENG, moi PDF mot gia tri (probe dau tien
% dat ca 4 gia tri trong cung mot file va tach doan bang khoang cach doc -> ham tach
% gop ca 4 doan lipsum thanh 1, trung vi ra 15.54pt la TRON cua 4 gia tri, vo nghia;
% \typeout khong xuat hien trong PDF nen khong dung lam moc duoc).
%   baselineskip tu nhien cua TeX Gyre Pagella 10pt = 11.954pt
%   (7 gia tri cho cung mot he so: 11.95-11.96pt, sai so 0.01pt)
%   1.15 -> 13.75pt   1.25 -> 14.94pt   1.29 -> 15.42pt   1.30 -> 15.54pt
% 2 bai HUJOS-TT da xuat ban do duoc pitch 15.48pt (n=225 va n=161 khoang cach).
% => \linespread{1.29} cho 15.42pt, lech target -0.06pt. Chon 1.29.
% Dung 1.15 (13.75pt) se chat hon bai cua chinh toa soan 1.13 lan.
\linespread{1.29}
\setlength{\parskip}{3pt plus 0.5pt}
\setlength{\parindent}{1.0cm}

% \raggedbottom la BAT BUOC o day. \documentclass[10pt,twoside] mac dinh bat
% \flushbottom, tuc LaTeX keo gian KEO DOC de moi trang cao bang nhau. Hau qua do
% duoc (experiments/verify_layout.py): parskip that cua ban em do ra +7.97pt thay vi
% 3pt da dat, va khoang cach tai ranh doan phan tan 16.9/21.7/23.9pt trong khi 2 bai
% da dang gom chat 18.1/21.5pt. \raggedbottom de trang ket thuc tu nhien, khong keo.
% parskip GIU 3pt: 2 bai da dang do duoc +2.64pt tai ranh doan (ca hai bai giong
% het nhau), khop "3 points after paragraph spacing" cua quy dinh van ban.
\raggedbottom

% ---- toan hoc ----
% KHONG dung \numberwithin{equation}{section}: quy dinh HUJOS-TT yeu cau cong thuc
% danh so lien tuc va duoc goi la "Eq. (1)", nen danh so theo section se ra (1.1)
% va lam sai ca cach goi lan thu tu trong van ban.
%
% KHONG nap amssymb CUNG unicode-math. Da do bang build that: hai goi trung dinh
% nghia 4 lenh va dung bien dich voi 4 loi fatal
%   ! LaTeX Error: Command `\eth' already defined.
%   ! LaTeX Error: Command `\smallsetminus' already defined.
%   ! LaTeX Error: Command `\digamma' already defined.
%   ! LaTeX Error: Command `\backepsilon' already defined.
% unicode-math (nap o tren, kem TeX Gyre Pagella Math) da cung cap day du symbol
% cua amssymb, nen bo amssymb la cach sua dung thay vi \let ghi de tung lenh.
\usepackage{amsmath}

% ---- bang: booktabs + caption tren bang ----
\usepackage{booktabs}
% Caption canh GIUA: do tu 2 bai da xuat ban (experiments/measure_venue_style.py),
% 17/17 caption 'Table n.'/'Fig. n.' co mid = 9.50cm dung tam vung chu
% (2.00..17.00cm), ke ca caption dai 2 dong (x=[2.79,16.21] va x=[2.74,16.09]).
% Ban cu dung justification=raggedright + singlelinecheck=false = canh TRAI,
% khong khop bai nao. font=small (9pt) giu nguyen: do duoc 9.14pt o bai da dang.
% labelfont=bf giu nguyen: quy dinh noi "Figure numbers and table numbers are
% printed in bold". labelsep=period giu nguyen: bai dang that la "Fig. 1." / "Table 1."
\usepackage[font=small,labelfont=bf,labelsep=period,
            skip=4pt,justification=centering]{caption}
\captionsetup[table]{position=top}
\captionsetup[figure]{position=bottom}
% Ban tieng Viet: nhan float phai la "Bang 1." / "Hinh 1." chu khong phai "Table 1."
% Quy dinh HUJOS ghi "table titles ABOVE the table, figure titles BELOW the figure"
% va "table/figure numbers printed in BOLD" -- vi tri va do dam khong doi, chi doi
% CHU cua nhan.
\ifhujosvi
  \captionsetup[table]{name=Bảng}
  \captionsetup[figure]{name=Hình}
\fi

% ---- heading theo so do tu template ----
\usepackage{titlesec}
% Bai da dang: so muc "1" nam o x=[2.00,2.33]cm, CHU "Introduction" bat dau o
% x=3.00cm (PyMuPDF tach thanh 2 span cung dong y=19.3). Tuc khoang cach tu CUOI
% nhan den dau chu = 0.79cm (do lai lan 2: het nhan 2.21 -> chu 3.00, ca 2 bai
% deu ra gap 0.79cm). Ban cu dung 0.5em (~0.21cm o 12pt) -> chu nam o
% ~2.25cm; ban sua 0.67cm cho 2.88cm, van lech 0.12cm -> chot 0.79cm.
% Subsection GIU 0.5em: bai #2 do duoc "2.1 Object Detection from UAV Imagery"
% bat dau o x0=2.00cm, khong thut them.
\titleformat{\section}{\normalfont\bfseries\fontsize{12}{14}\selectfont}
  {\thesection}{0.79cm}{}
\titlespacing*{\section}{0pt}{18pt}{12pt}
\titleformat{\subsection}{\normalfont\bfseries\fontsize{10}{12}\selectfont}
  {\thesubsection}{0.5em}{}
\titlespacing*{\subsection}{0pt}{18pt}{8pt}
\titleformat{\subsubsection}[runin]{\normalfont\itshape}
  {\thesubsubsection}{0.5em}{}[.]
\titlespacing*{\subsubsection}{0pt}{12pt}{0.5em}
% 18 pt space after text in each section (template yeu cau)
\titleformat{name=\section,numberless}{\normalfont\bfseries\fontsize{12}{14}\selectfont}
  {}{0pt}{}
\titlespacing*{name=\section,numberless}{0pt}{18pt}{12pt}

% ---- dinh ly / chung minh ----
\usepackage{amsthm}
% Ban tieng Viet doi ten moi truong dinh ly (KHONG doi noi dung dinh ly).
% Ten tieng Viet dung theo thong le toan hoc Viet Nam: proposition = Menh de.
\ifhujosvi
  \newtheorem{proposition}{Mệnh đề}
  \newtheorem{remark}{Ghi chú}
  \renewcommand{\proofname}{Chứng minh}
\else
  \newtheorem{proposition}{Proposition}
  \newtheorem{remark}{Remark}
  \renewcommand{\proofname}{Proof}
\fi

% ---- trích dẫn Vancouver (ngoặc vuông) từ 2026 ----
\usepackage[numbers,sort&compress]{natbib}
\bibpunct{[}{]}{,}{n}{}{,}

% ---- link (an vien, khong doi mau) ----
% PHAI \PassOptionsToPackage TRUOC hyperref: hyperref tu nap url ben trong, nen
% viet \usepackage[hyphens]{url} o SAU se dung "Option clash for package url"
% va dung bien dich. Da tra gia: build bao P3=1 ma van xuat PDF 13 trang, trong
% do 3 overfull cua .bbl giu nguyen 189.81/139.74/118.98pt -> hyphens chua an.
\PassOptionsToPackage{hyphens}{url}
\usepackage[hidelinks]{hyperref}

% ---- header/footer nhu template: khong tu them so trang ngoai cac thu da set ----
\usepackage{fancyhdr}
\pagestyle{fancy}
\fancyhf{}
\renewcommand{\headrulewidth}{0pt}
% Header: ban tieng Viet ghi ten tac gia khong dau (dung quy uoc cua chinh tap chi
% o cac bai tieng Viet da dang) va ten tap chi tieng Viet.
\ifhujosvi
  \fancyhead[LE]{\footnotesize\itshape M. X. Văn và cs.}
  \fancyhead[RO]{\footnotesize\itshape Tạp chí Khoa học Đại học Huế: Kỹ thuật và Công nghệ}
\else
  \fancyhead[LE]{\footnotesize\itshape X. V. Mai et al.}
  \fancyhead[RO]{\footnotesize\itshape Hue University Journal of Science: Techniques and Technology}
\fi
\fancyfoot[C]{\footnotesize\thepage}

% ---- phu luc: doi ten "Appendix" thanh "Appendix" (giu nguyen) ----
\usepackage{appendix}

% ---- tien ich ----
% [hyphens] la BAT BUOC: ba overfull nang nhat cua bai (189.81 / 139.74 /
% 118.98 pt, tuc chu loi ra ngoai trang) deu nam trong .bbl va deu do URL dai
% khong ngat duoc. Mac dinh goi url chi ngat o dau '/' va '.', ma ten mien nhu
% 'technical-explainer-carbon-credits-for-improved-rice-cultivation' khong co '/'.
% Option duoc truyen bang \PassOptionsToPackage o tren (truoc hyperref) de tranh
% option clash, nen o day nap khong kem option.
\usepackage{url}
\usepackage{xspace}
\usepackage{microtype}
\usepackage{array}

% Moi truong liet ke ky hieu cong thuc ('trong do:'). Dinh nghia o day chu khong o
% blocks/ vi cac lenh \setlength chua CHU SO (0.5pt, 0pt) se bi gate
% check_vi_numbers dem nhu so lieu nghien cuu -> lam loang gate.
\newenvironment{symbolist}
  {\begin{description}\footnotesize\setlength{\itemsep}{0.5pt}\setlength{\parsep}{0pt}}
  {\end{description}}


% graphicx la BAT BUOC cho \resizebox. Thieu no thi 24 loi fatal
% "Undefined control sequence" + "Missing number" + "Illegal unit of measure"
% (vi \resizebox bi phan tich sai). Da tra gia: ban IEEE co goi nay, ban HUJOS
% quen -> build ra PDF 13 trang nhung la RAC, khong duoc bao cao la thanh cong.
\usepackage{graphicx}

% \fittocol: chi THU NHO bang khi no that su rong hon cot, KHONG BAO GIO phong to.
%
% Tai sao khong dung \sbox de do: cot p{} chua \par, ma \sbox = \savebox trong
% \hbox nen khong do duoc vat lieu dang doan van. Do bang \sbox cho ra 7.92-8.01cm
% trong khi build that bao overfull 18-26pt (~8.8cm) -> so do sai, dan den ket luan
% sai rang "bang khong tran". lrbox la cong cu dung cho hop chua \par.
%
% Tai sao khong dung \resizebox{\columnwidth}{!} truc tiep: no CUONG BUC dung bang
% kho cot, nen bang hep hon (vd tab_baseline 13.69cm trong table* 16.5cm) se bi
% PHONG TO ~20%, chu to hon thiet ke.
\newsavebox{\fitbox}
\newcommand{\fittocol}[1]{%
  \begin{lrbox}{\fitbox}#1\end{lrbox}%
  \ifdim\wd\fitbox>\columnwidth
    \resizebox{\columnwidth}{!}{\usebox{\fitbox}}%
  \else
    \usebox{\fitbox}%
  \fi}

% Cho phep keo gian khoang trang them 2em khi khong tim duoc diem ngat.
% Ly do: cac doan chua nhieu \ref lien tiep (vd
% "(Sec.~\ref{sec:umax}, Sec.~\ref{sec:v2reach}, Table~\ref{tab:umax},
% Table~\ref{tab:f2reach})") tao ra chuoi khong ngat duoc, gay overfull
% 34.89pt va 26.16pt trong 01_intro.tex. Day la cach sua khong doi noi dung.
\setlength{\emergencystretch}{2em}

% Kho cot that: MOT cot, vung chu 19 - 2 - 2 = 15.00 cm
% (bai da xuat ban do duoc 15.01 cm — khop). \columnwidth = \textwidth o day.
% Day la con so ma experiments/make_tables_hujos.py dung de gate chieu rong bang.
\newlength{\HUJOScol}
\setlength{\HUJOScol}{15.0cm}

% ---------------- FRONT MATTER: moi con so DO TU 2 BAI HUJOS-TT DA XUAT BAN --------
% experiments/measure_fonts.py (PyMuPDF, doc theo SPAN nen biet dung font cua tung
% manh chu, khong gop dong):
%   title        PalatinoLinotype-Bold 14.04pt, canh giua       -> em: 14pt bold center  (khop)
%   authors      canh GIUA: mid = 9.59 va 9.53 cm (tam vung chu 9.5)
%   affiliations canh GIUA: mid = 9.54 cm (ca 4 dong cua bai #1)
%   abstract     9.00pt, x = [3.00, 16.08] cm -> THUT VAO 1.00 cm ca hai ben,
%                trong khi than bai x = [2.00, 17.08] cm (khong thut)
%   keywords     cung nam trong khoi thut 1.00 cm (x bat dau 3.00)
%   caption      9.00pt, chi NHAN in dam: span 'Fig. 1.' = PalatinoLinotype-Bold,
%                span con lai = PalatinoLinotype-Roman (10/10 caption bai #1,
%                8/8 bai #2) -> labelfont=bf + font=small cua em da khop san
% Ban cu de \raggedright cho authors/affil va khong thut le abstract: lech ca hai
% bai da dang o dung nhung cho nay.
\newcommand{\hujostitle}[1]{%
  {\centering\normalfont\bfseries\fontsize{14}{17}\selectfont #1\par}}
\newcommand{\hujosauthors}[1]{%
  {\centering\normalfont\bfseries\fontsize{10}{12}\selectfont #1\par}}
\newcommand{\hujosaffil}[1]{%
  {\centering\normalfont\fontsize{9}{11}\selectfont #1\par}}
% \leftskip/\rightskip = 1cm de abstract + keywords thut vao dung nhu bai da dang.
% Keywords nam BEN TRONG hujosabstract trong main_hujos.tex nen cung duoc thut.
% CO CHU ABSTRACT = 9pt, KHONG PHAI 11pt.
% LOI EM VUA GAY RA TRONG CHINH LAN SUA NAY: go \fontsize{11}{13} thay vi {9}{11}
% -> abstract se to hon than bai (10pt), nguoc hoan toan voi so do da do
% (measure_fonts.py: abstract 9.00pt o ca 2 bai; body 9.96pt). Da sua lai 9/11.
\newenvironment{hujosabstract}
  {\par\vspace{10pt}\normalfont\fontsize{9}{11}\selectfont
   \leftskip=1cm\rightskip=1cm\parindent=0pt\noindent\ignorespaces}
  {\par}
. Ban tieng Anh
% khong dinh nghia gi => \ifhujosvi = false => MOI thu giong het truoc khi sua
% (da kiem chung: main_hujos.pdf van 14 trang, 0 undefined, 0 overfull).
% Ly do dung macro \HUJOSVI thay vi babel/polyglossia: tieng Viet trong TeX Live
% chi co `zerohyph.tex` (khong co mau ngat am), nen nap polyglossia chi them
% rui ro xung dot voi unicode-math ma khong duoc loi ich nao. Tu tieng Viet ngan
% va tach bang dau cach nen diem ngat dong rat day du.
\newif\ifhujosvi
\ifdefined\HUJOSVI\hujosvitrue\fi

% ---- hinh hoc ----
% DA XAC MINH BANG BAI DA XUAT BAN THAT, khong con dua vao template nao.
%
% LOI DA MAC (ghi lai de khong tai pham): ban dau lay hinh hoc tu
% /tmp/hujos_tpl.docx = A4 20.5x29.7cm, le 1.5cm tren/duoi, HAI cot, columnsep
% 425 twips. File do CO that va cac con so do DUNG (425 twips nam trong the con
% <w:col w:space="425" w:w="4465.5">, khong phai trong <w:cols/>). Nhung HEADER
% BEN TRONG cua no ghi "Hue University Journal of Science: Natural Science
% Vol. 129, No. 1B" — do la template cua HUJOS KHOA HOC TU NHIEN, mot tap chi
% KHAC trong cung he thong. Bai nay nop HUJOS-TT (Ky thuat & Cong nghe).
% Hau qua: dung layout cua tap chi khac, vung chu 16.74cm thay vi 15.01cm,
% rong hon 11.5% -> so trang do duoc (13) khong co gia tri.
%
% BANG CHUNG QUYET DINH = 2 bai HUJOS-TT da xuat ban, tai tu
% citation_pdf_url cua chinh trang bai (article 7830, 7862):
%   pdfinfo     : 538.7 x 765.5 pt = 19.00 x 27.00 cm  (ca hai bai)
%   vung chu    : xMin 2.00cm -> xMax 17.01cm, rong 15.01 cm
%   le          : trai 2.00, phai 1.99 cm  -> le 2cm deu, dung nhu quy dinh van ban
%   SO COT      : chu bat dau trong dai giua trang = 20.1% va 20.5% (day chu)
%                 -> MOT COT. (Ban hai cot chi co 4.9%, vi ranh cot trong.)
%   font        : PalatinoLinotype-Roman/Bold/Italic + Cambria + Arial-BoldMT
% Kiem cheo: Author Guidelines cua HUJOS-TT noi "19 cm x 27 cm ... margins on
% all sides of 2 cm ... Palatino Linotype ... font size 10", va NHAC DEN
% "column" 0 LAN. Ba nguon (quy dinh van ban + 2 bai da xuat ban + template docx
% ref_hujos*.docx 19x27 le 2 1 cot) nay THONG NHAT. Chi co hujos_tpl.docx lech,
% va no thuoc tap chi khac.
\usepackage[paperwidth=19cm,paperheight=27cm,
            top=2cm,bottom=2cm,left=2cm,right=2cm,
            headsep=0.4cm,footskip=0.6cm]{geometry}

% ---- font: Palatino (TeX Gyre Pagella) + math khop ----
\usepackage{fontspec}
\setmainfont{TeX Gyre Pagella}
\usepackage{unicode-math}
\setmathfont{TeX Gyre Pagella Math}

% ---- MOT COT (xac nhan bang 2 bai HUJOS-TT da xuat ban) ----
% \columnsep khong con tac dung khi document khong \twocolumn; giu gia tri nay
% chi de neu ai do build lai 2 cot thi khong vo tinh dung khoang cach cua tap chi
% khac. Vung chu that = 15.01 cm theo bai da xuat ban.
\setlength{\columnsep}{0.75cm}

% ---- gian dong + indent doan: DO TU 2 BAI HUJOS-TT DA XUAT BAN ----
% experiments/measure_indent.py (loc dung bang tieu chi "doan truoc ket thuc som
% hon le phai 0.5cm", khong dem moi dong thut le vi nham voi dong trong bang):
%   indent dau doan = 1.00 cm  (bai #1: 31/72 dong; bai #2: 30/58 dong — trung vi
%                               cua ca hai dung bang 1.00)
%   khoang cach doan = pitch tai cho bat dau doan tru them +3.01pt (bai #1)
%                                            va +6.00pt (bai #2) so voi pitch thuong
% Quy dinh van ban: "6 points before and 3 points after paragraph spacing".
% template docx docDefaults: w:spacing before="120" after="60" line="276"
% (= 6pt / 3pt / 1.15). Hai nguon nay khop nhau -> chon parskip 3pt.
%
% Ban cu de parindent 0.45cm va parskip 0pt: KHONG khop nguon nao.
%
% "1.15" CUA WORD KHAC \linespread{1.15} CUA LATEX — da giai bang phep do, khong
% con la xung dot phai doan. Word nhan 1.15 voi (ascent+descent+lineGap) cua FONT;
% LaTeX \linespread nhan voi baselineskip cua CLASS.
%
% experiments/probe_linespread.py: 7 PDF RIENG, moi PDF mot gia tri (probe dau tien
% dat ca 4 gia tri trong cung mot file va tach doan bang khoang cach doc -> ham tach
% gop ca 4 doan lipsum thanh 1, trung vi ra 15.54pt la TRON cua 4 gia tri, vo nghia;
% \typeout khong xuat hien trong PDF nen khong dung lam moc duoc).
%   baselineskip tu nhien cua TeX Gyre Pagella 10pt = 11.954pt
%   (7 gia tri cho cung mot he so: 11.95-11.96pt, sai so 0.01pt)
%   1.15 -> 13.75pt   1.25 -> 14.94pt   1.29 -> 15.42pt   1.30 -> 15.54pt
% 2 bai HUJOS-TT da xuat ban do duoc pitch 15.48pt (n=225 va n=161 khoang cach).
% => \linespread{1.29} cho 15.42pt, lech target -0.06pt. Chon 1.29.
% Dung 1.15 (13.75pt) se chat hon bai cua chinh toa soan 1.13 lan.
\linespread{1.29}
\setlength{\parskip}{3pt plus 0.5pt}
\setlength{\parindent}{1.0cm}

% \raggedbottom la BAT BUOC o day. \documentclass[10pt,twoside] mac dinh bat
% \flushbottom, tuc LaTeX keo gian KEO DOC de moi trang cao bang nhau. Hau qua do
% duoc (experiments/verify_layout.py): parskip that cua ban em do ra +7.97pt thay vi
% 3pt da dat, va khoang cach tai ranh doan phan tan 16.9/21.7/23.9pt trong khi 2 bai
% da dang gom chat 18.1/21.5pt. \raggedbottom de trang ket thuc tu nhien, khong keo.
% parskip GIU 3pt: 2 bai da dang do duoc +2.64pt tai ranh doan (ca hai bai giong
% het nhau), khop "3 points after paragraph spacing" cua quy dinh van ban.
\raggedbottom

% ---- toan hoc ----
% KHONG dung \numberwithin{equation}{section}: quy dinh HUJOS-TT yeu cau cong thuc
% danh so lien tuc va duoc goi la "Eq. (1)", nen danh so theo section se ra (1.1)
% va lam sai ca cach goi lan thu tu trong van ban.
%
% KHONG nap amssymb CUNG unicode-math. Da do bang build that: hai goi trung dinh
% nghia 4 lenh va dung bien dich voi 4 loi fatal
%   ! LaTeX Error: Command `\eth' already defined.
%   ! LaTeX Error: Command `\smallsetminus' already defined.
%   ! LaTeX Error: Command `\digamma' already defined.
%   ! LaTeX Error: Command `\backepsilon' already defined.
% unicode-math (nap o tren, kem TeX Gyre Pagella Math) da cung cap day du symbol
% cua amssymb, nen bo amssymb la cach sua dung thay vi \let ghi de tung lenh.
\usepackage{amsmath}

% ---- bang: booktabs + caption tren bang ----
\usepackage{booktabs}
% Caption canh GIUA: do tu 2 bai da xuat ban (experiments/measure_venue_style.py),
% 17/17 caption 'Table n.'/'Fig. n.' co mid = 9.50cm dung tam vung chu
% (2.00..17.00cm), ke ca caption dai 2 dong (x=[2.79,16.21] va x=[2.74,16.09]).
% Ban cu dung justification=raggedright + singlelinecheck=false = canh TRAI,
% khong khop bai nao. font=small (9pt) giu nguyen: do duoc 9.14pt o bai da dang.
% labelfont=bf giu nguyen: quy dinh noi "Figure numbers and table numbers are
% printed in bold". labelsep=period giu nguyen: bai dang that la "Fig. 1." / "Table 1."
\usepackage[font=small,labelfont=bf,labelsep=period,
            skip=4pt,justification=centering]{caption}
\captionsetup[table]{position=top}
\captionsetup[figure]{position=bottom}
% Ban tieng Viet: nhan float phai la "Bang 1." / "Hinh 1." chu khong phai "Table 1."
% Quy dinh HUJOS ghi "table titles ABOVE the table, figure titles BELOW the figure"
% va "table/figure numbers printed in BOLD" -- vi tri va do dam khong doi, chi doi
% CHU cua nhan.
\ifhujosvi
  \captionsetup[table]{name=Bảng}
  \captionsetup[figure]{name=Hình}
\fi

% ---- heading theo so do tu template ----
\usepackage{titlesec}
% Bai da dang: so muc "1" nam o x=[2.00,2.33]cm, CHU "Introduction" bat dau o
% x=3.00cm (PyMuPDF tach thanh 2 span cung dong y=19.3). Tuc khoang cach tu CUOI
% nhan den dau chu = 0.79cm (do lai lan 2: het nhan 2.21 -> chu 3.00, ca 2 bai
% deu ra gap 0.79cm). Ban cu dung 0.5em (~0.21cm o 12pt) -> chu nam o
% ~2.25cm; ban sua 0.67cm cho 2.88cm, van lech 0.12cm -> chot 0.79cm.
% Subsection GIU 0.5em: bai #2 do duoc "2.1 Object Detection from UAV Imagery"
% bat dau o x0=2.00cm, khong thut them.
\titleformat{\section}{\normalfont\bfseries\fontsize{12}{14}\selectfont}
  {\thesection}{0.79cm}{}
\titlespacing*{\section}{0pt}{18pt}{12pt}
\titleformat{\subsection}{\normalfont\bfseries\fontsize{10}{12}\selectfont}
  {\thesubsection}{0.5em}{}
\titlespacing*{\subsection}{0pt}{18pt}{8pt}
\titleformat{\subsubsection}[runin]{\normalfont\itshape}
  {\thesubsubsection}{0.5em}{}[.]
\titlespacing*{\subsubsection}{0pt}{12pt}{0.5em}
% 18 pt space after text in each section (template yeu cau)
\titleformat{name=\section,numberless}{\normalfont\bfseries\fontsize{12}{14}\selectfont}
  {}{0pt}{}
\titlespacing*{name=\section,numberless}{0pt}{18pt}{12pt}

% ---- dinh ly / chung minh ----
\usepackage{amsthm}
% Ban tieng Viet doi ten moi truong dinh ly (KHONG doi noi dung dinh ly).
% Ten tieng Viet dung theo thong le toan hoc Viet Nam: proposition = Menh de.
\ifhujosvi
  \newtheorem{proposition}{Mệnh đề}
  \newtheorem{remark}{Ghi chú}
  \renewcommand{\proofname}{Chứng minh}
\else
  \newtheorem{proposition}{Proposition}
  \newtheorem{remark}{Remark}
  \renewcommand{\proofname}{Proof}
\fi

% ---- trích dẫn Vancouver (ngoặc vuông) từ 2026 ----
\usepackage[numbers,sort&compress]{natbib}
\bibpunct{[}{]}{,}{n}{}{,}

% ---- link (an vien, khong doi mau) ----
% PHAI \PassOptionsToPackage TRUOC hyperref: hyperref tu nap url ben trong, nen
% viet \usepackage[hyphens]{url} o SAU se dung "Option clash for package url"
% va dung bien dich. Da tra gia: build bao P3=1 ma van xuat PDF 13 trang, trong
% do 3 overfull cua .bbl giu nguyen 189.81/139.74/118.98pt -> hyphens chua an.
\PassOptionsToPackage{hyphens}{url}
\usepackage[hidelinks]{hyperref}

% ---- header/footer nhu template: khong tu them so trang ngoai cac thu da set ----
\usepackage{fancyhdr}
\pagestyle{fancy}
\fancyhf{}
\renewcommand{\headrulewidth}{0pt}
% Header: ban tieng Viet ghi ten tac gia khong dau (dung quy uoc cua chinh tap chi
% o cac bai tieng Viet da dang) va ten tap chi tieng Viet.
\ifhujosvi
  \fancyhead[LE]{\footnotesize\itshape M. X. Văn và cs.}
  \fancyhead[RO]{\footnotesize\itshape Tạp chí Khoa học Đại học Huế: Kỹ thuật và Công nghệ}
\else
  \fancyhead[LE]{\footnotesize\itshape X. V. Mai et al.}
  \fancyhead[RO]{\footnotesize\itshape Hue University Journal of Science: Techniques and Technology}
\fi
\fancyfoot[C]{\footnotesize\thepage}

% ---- phu luc: doi ten "Appendix" thanh "Appendix" (giu nguyen) ----
\usepackage{appendix}

% ---- tien ich ----
% [hyphens] la BAT BUOC: ba overfull nang nhat cua bai (189.81 / 139.74 /
% 118.98 pt, tuc chu loi ra ngoai trang) deu nam trong .bbl va deu do URL dai
% khong ngat duoc. Mac dinh goi url chi ngat o dau '/' va '.', ma ten mien nhu
% 'technical-explainer-carbon-credits-for-improved-rice-cultivation' khong co '/'.
% Option duoc truyen bang \PassOptionsToPackage o tren (truoc hyperref) de tranh
% option clash, nen o day nap khong kem option.
\usepackage{url}
\usepackage{xspace}
\usepackage{microtype}
\usepackage{array}

% Moi truong liet ke ky hieu cong thuc ('trong do:'). Dinh nghia o day chu khong o
% blocks/ vi cac lenh \setlength chua CHU SO (0.5pt, 0pt) se bi gate
% check_vi_numbers dem nhu so lieu nghien cuu -> lam loang gate.
\newenvironment{symbolist}
  {\begin{description}\footnotesize\setlength{\itemsep}{0.5pt}\setlength{\parsep}{0pt}}
  {\end{description}}


% graphicx la BAT BUOC cho \resizebox. Thieu no thi 24 loi fatal
% "Undefined control sequence" + "Missing number" + "Illegal unit of measure"
% (vi \resizebox bi phan tich sai). Da tra gia: ban IEEE co goi nay, ban HUJOS
% quen -> build ra PDF 13 trang nhung la RAC, khong duoc bao cao la thanh cong.
\usepackage{graphicx}

% \fittocol: chi THU NHO bang khi no that su rong hon cot, KHONG BAO GIO phong to.
%
% Tai sao khong dung \sbox de do: cot p{} chua \par, ma \sbox = \savebox trong
% \hbox nen khong do duoc vat lieu dang doan van. Do bang \sbox cho ra 7.92-8.01cm
% trong khi build that bao overfull 18-26pt (~8.8cm) -> so do sai, dan den ket luan
% sai rang "bang khong tran". lrbox la cong cu dung cho hop chua \par.
%
% Tai sao khong dung \resizebox{\columnwidth}{!} truc tiep: no CUONG BUC dung bang
% kho cot, nen bang hep hon (vd tab_baseline 13.69cm trong table* 16.5cm) se bi
% PHONG TO ~20%, chu to hon thiet ke.
\newsavebox{\fitbox}
\newcommand{\fittocol}[1]{%
  \begin{lrbox}{\fitbox}#1\end{lrbox}%
  \ifdim\wd\fitbox>\columnwidth
    \resizebox{\columnwidth}{!}{\usebox{\fitbox}}%
  \else
    \usebox{\fitbox}%
  \fi}

% Cho phep keo gian khoang trang them 2em khi khong tim duoc diem ngat.
% Ly do: cac doan chua nhieu \ref lien tiep (vd
% "(Sec.~\ref{sec:umax}, Sec.~\ref{sec:v2reach}, Table~\ref{tab:umax},
% Table~\ref{tab:f2reach})") tao ra chuoi khong ngat duoc, gay overfull
% 34.89pt va 26.16pt trong 01_intro.tex. Day la cach sua khong doi noi dung.
\setlength{\emergencystretch}{2em}

% Kho cot that: MOT cot, vung chu 19 - 2 - 2 = 15.00 cm
% (bai da xuat ban do duoc 15.01 cm — khop). \columnwidth = \textwidth o day.
% Day la con so ma experiments/make_tables_hujos.py dung de gate chieu rong bang.
\newlength{\HUJOScol}
\setlength{\HUJOScol}{15.0cm}

% ---------------- FRONT MATTER: moi con so DO TU 2 BAI HUJOS-TT DA XUAT BAN --------
% experiments/measure_fonts.py (PyMuPDF, doc theo SPAN nen biet dung font cua tung
% manh chu, khong gop dong):
%   title        PalatinoLinotype-Bold 14.04pt, canh giua       -> em: 14pt bold center  (khop)
%   authors      canh GIUA: mid = 9.59 va 9.53 cm (tam vung chu 9.5)
%   affiliations canh GIUA: mid = 9.54 cm (ca 4 dong cua bai #1)
%   abstract     9.00pt, x = [3.00, 16.08] cm -> THUT VAO 1.00 cm ca hai ben,
%                trong khi than bai x = [2.00, 17.08] cm (khong thut)
%   keywords     cung nam trong khoi thut 1.00 cm (x bat dau 3.00)
%   caption      9.00pt, chi NHAN in dam: span 'Fig. 1.' = PalatinoLinotype-Bold,
%                span con lai = PalatinoLinotype-Roman (10/10 caption bai #1,
%                8/8 bai #2) -> labelfont=bf + font=small cua em da khop san
% Ban cu de \raggedright cho authors/affil va khong thut le abstract: lech ca hai
% bai da dang o dung nhung cho nay.
\newcommand{\hujostitle}[1]{%
  {\centering\normalfont\bfseries\fontsize{14}{17}\selectfont #1\par}}
\newcommand{\hujosauthors}[1]{%
  {\centering\normalfont\bfseries\fontsize{10}{12}\selectfont #1\par}}
\newcommand{\hujosaffil}[1]{%
  {\centering\normalfont\fontsize{9}{11}\selectfont #1\par}}
% \leftskip/\rightskip = 1cm de abstract + keywords thut vao dung nhu bai da dang.
% Keywords nam BEN TRONG hujosabstract trong main_hujos.tex nen cung duoc thut.
% CO CHU ABSTRACT = 9pt, KHONG PHAI 11pt.
% LOI EM VUA GAY RA TRONG CHINH LAN SUA NAY: go \fontsize{11}{13} thay vi {9}{11}
% -> abstract se to hon than bai (10pt), nguoc hoan toan voi so do da do
% (measure_fonts.py: abstract 9.00pt o ca 2 bai; body 9.96pt). Da sua lai 9/11.
\newenvironment{hujosabstract}
  {\par\vspace{10pt}\normalfont\fontsize{9}{11}\selectfont
   \leftskip=1cm\rightskip=1cm\parindent=0pt\noindent\ignorespaces}
  {\par}


\begin{document}

% Bang cua phu luc danh so S1..S5 de khop nhan 'Bang SN' ma bai chinh trich dan.
\renewcommand{\thetable}{S\arabic{table}}

\begingroup
  \hujostitle{Phụ lục\\
  Điểm mù cố hữu của chứng thực vật lý đối với nhật ký tưới và giải pháp
  neo theo thiết bị trong đo đạc, báo cáo, thẩm định lúa phát thải thấp}

  \vspace{8pt}
  \hujosauthors{Mai Xuân Văn$^{1}$, Lê Đức Minh Phương$^{2,*}$, Nguyễn Đăng Trị$^{3}$,
  Trương Khánh Duy$^{2}$, Nguyễn Tương Tri$^{2,4,*}$}

  \vspace{4pt}
  \hujosaffil{$^{1}$ Trường Trung học phổ thông (THPT) Thuận Hóa, Trường Đại học Sư phạm, Đại học Huế, Huế, Việt Nam\\
  $^{2}$ Viện Đào tạo Mở và Công nghệ thông tin, Đại học Huế, Huế, Việt Nam\\
  $^{3}$ Trường Đại học Gia Định, Thành phố Hồ Chí Minh, Việt Nam\\
  $^{4}$ Trường Đại học Sư phạm, Đại học Huế, 34 Lê Lợi, Huế, Việt Nam}

  \vspace{4pt}
  \hujosaffil{$^{*}$ Đồng tác giả liên hệ: Lê Đức Minh Phương. E-mail: Ldmphuong@hueuni.edu.vn;\\
  Nguyễn Tương Tri. E-mail: ntuongtri@hueuni.edu.vn}
\endgroup

\vspace{10pt}

% Stub nhan: TU SINH tu main_hujos_vi.aux. Build lai bai chinh roi chay
%   python3 experiments/make_supp_stubs_vi.py
% neu bai chinh doi danh so cong thuc/muc/bang.
\InputIfFileExists{supp_stubs_vi.tex}{}{%
  \errmessage{THIEU supp_stubs_vi.tex: chay python3 experiments/make_supp_stubs_vi.py}}

% ---------------------------------------------------------------- S1
\section*{S1. Các lớp gian lận và mỗi tầng phòng thủ quan sát được gì}

\begin{table}[!h]
\centering
\caption{Ba đại lượng làm giả được của \eqref{eq:three} trong bài chính và lớp gian lận
nào làm thay đổi từng đại lượng. ``đúng'' nghĩa là đại lượng đó được giữ nguyên, nên tầng
phòng thủ tương ứng bị mù theo Mệnh đề~\ref{prop:blind} của bài chính.}
\footnotesize
% Sinh TU DONG tu tables/tab_threat.tex (ban tieng Anh, von sinh tu outputs/*.csv)
% boi experiments/make_tables_hujos_vi.py. KHONG sua tay.
% Nhan dich qua bang tra cuong che; so chi doi dinh dang sang kieu Viet Nam (dau phay thap phan).
% Da assert: cung so dong, cung so vach, cung tap chu so voi ban EN.

% Sinh TU DONG tu outputs/eval_theta_*.csv boi experiments/make_tables_anchored.py
% KHONG sua tay: moi con so deu truy nguon duoc ve CSV.

\fittocol{%
\begin{tabular}{@{}>{\raggedright\arraybackslash}p{2.85cm}>{\raggedright\arraybackslash}p{0.95cm}>{\raggedright\arraybackslash}p{1.05cm}>{\raggedright\arraybackslash}p{0.95cm}>{\raggedright\arraybackslash}p{2.20cm}@{}}
\toprule
Biến thể gian lận & Tổng mm & Từng độ sâu & Thời điểm & Làm giả cái gì \\
\midrule
Nông dân trung thực (điều kiện ruộng) & đúng & đúng & đúng & Mốc: bảy trục khác biệt thực tế đang bật \\
Thêm lệnh tưới & SAI & đúng & đúng & Thổi phồng hồ sơ hoạt động quản lý \\
Độ sâu khai không khả thi & SAI & SAI & đúng & Một phần ba số mục đặt bằng toàn bộ dung tích chứa \\
Bỏ sót lệnh thật & SAI & đúng & đúng & Xoá cửa sổ 96 h để tạo ra một pha khô \\
Gộp khai báo, 24 h & đúng & thay đổi & SAI & Cùng lượng nước, ít lần tưới hơn và lớn hơn \\
Gộp khai báo, 96 h & đúng & thay đổi & SAI & Cùng lượng nước, các lần tưới bám vào lưới thô \\
Dịch thời điểm (duy lý) & đúng & đúng & SAI & Leo đồi theo từng giờ để tối đa số pha khô được cấp tín chỉ \\
Bên gian lận duy lý & đúng & đúng & SAI & Chọn cửa sổ gộp làm tối đa số pha khô \\
\bottomrule
\end{tabular}
}

\end{table}

% ---------------------------------------------------------------- S2
\section*{S2. Biên độ sâu chấp nhận được trên toàn $\Theta$}

\begin{table}[!h]
\centering
\caption{Biên độ sâu chấp nhận được $u_{\max}$ (\eqref{eq:umax} trong bài chính) trên
toàn dải của FAO-56 Bảng 2. Giá trị đơn lẻ $61$\,mm ban đầu chỉ là một dòng của bảng này.}
\footnotesize
% Sinh TU DONG tu tables/tab_umax.tex (ban tieng Anh, von sinh tu outputs/*.csv)
% boi experiments/make_tables_hujos_vi.py. KHONG sua tay.
% Nhan dich qua bang tra cuong che; so chi doi dinh dang sang kieu Viet Nam (dau phay thap phan).
% Da assert: cung so dong, cung so vach, cung tap chu so voi ban EN.

% Sinh TU DONG tu outputs/eval_theta_*.csv boi experiments/make_tables_anchored.py
% KHONG sua tay: moi con so deu truy nguon duoc ve CSV.

\fittocol{%
\begin{tabular}{@{}>{\raggedright\arraybackslash}p{3.00cm}>{\raggedright\arraybackslash}p{1.25cm}>{\raggedright\arraybackslash}p{1.85cm}>{\raggedright\arraybackslash}p{1.85cm}@{}}
\toprule
Loại đất (FAO-56 Bảng 2) & Biên độ (mm) & $u_{\max}$ tại $w_0=0$ & $u_{\max}$ tại $w_0=0{,}90$ \\
\midrule
đất thịt pha cát & 55,0 & 49,0 & 54,5 \\
đất thịt & 75,0 & 65,0 & 72,5 \\
đất thịt pha limon & 90,0 & 77,0 & 86,0 \\
đất thịt pha sét limon & 70,0 & 61,0 & 68,0 \\
đất sét & 80,0 & 69,0 & 77,0 \\
\midrule
\textbf{Dải phải công bố} & --- & \textbf{49,0} & \textbf{86,0} \\
\bottomrule
\end{tabular}
}

\end{table}

% ---------------------------------------------------------------- S3
\section*{S3. Khả năng kích hoạt của bộ lọc điểm héo}

\begin{table}[!h]
\centering
\caption{Khả năng kích hoạt của bộ lọc điểm héo trên $150$ cấu hình với nhật ký rỗng, tức
điều kiện thời tiết khô hạn nhất có được.}
\footnotesize
% Sinh TU DONG tu tables/tab_f2_unreachable.tex (ban tieng Anh, von sinh tu outputs/*.csv)
% boi experiments/make_tables_hujos_vi.py. KHONG sua tay.
% Nhan dich qua bang tra cuong che; so chi doi dinh dang sang kieu Viet Nam (dau phay thap phan).
% Da assert: cung so dong, cung so vach, cung tap chu so voi ban EN.

% Sinh TU DONG tu outputs/eval_theta_*.csv boi experiments/make_tables_anchored.py
% KHONG sua tay: moi con so deu truy nguon duoc ve CSV.

\fittocol{%
\begin{tabular}{@{}>{\raggedright\arraybackslash}p{1.55cm}c>{\raggedright\arraybackslash}p{1.75cm}>{\raggedright\arraybackslash}p{0.85cm}>{\raggedright\arraybackslash}p{2.60cm}@{}}
\toprule
Trạm & Năm & Cạn kiệt cắt ngọn lớn nhất & $V_2$ kích hoạt & Diễn giải \\
\midrule
Cần Thơ & 2024 & 0,998 & 0 & không bao giờ xuống dưới điểm héo \\
Cần Thơ & 2025 & 0,988 & 0 & không bao giờ xuống dưới điểm héo \\
Sóc Trăng & 2024 & 0,999 & 0 & không bao giờ xuống dưới điểm héo \\
Sóc Trăng & 2025 & 0,992 & 0 & không bao giờ xuống dưới điểm héo \\
Cà Mau & 2024 & 0,991 & 0 & không bao giờ xuống dưới điểm héo \\
Cà Mau & 2025 & 0,980 & 0 & không bao giờ xuống dưới điểm héo \\
\bottomrule
\end{tabular}
}

\end{table}

% ---------------------------------------------------------------- S4
% S4: ban TIENG VIET cua paper2/blocks/08_subsidence.tex.
% Chi duoc \input boi paper2_hujos/supplementary_vi.tex.
% NGUYEN TAC: giu nguyen moi con so va moi hedge; thap phan kieu Viet ({,});
% tham chieu cheo toi bai chinh dung \eqref/\ref (so lay tu main_hujos_vi.aux
% boi experiments/make_supp_stubs_vi.py) thay vi go tay "Eq.~(16)" — ban EN go tay
% nen da lech so voi ban VI (eq:umax EN=16, VI=15).
\section*{S4. Tính bất biến đối với sụt lún}

Tuân thủ AWD thường được phát biểu bằng một mực nước ngầm hoặc độ sâu ngập tuyệt
đối, ví dụ $15$\,cm dưới mặt ruộng, tức ngưỡng tưới lại của nghiệm thức Safe-AWD
trong~\cite{balaine2019}. Ở một đồng bằng đang sụt lún, cao độ mặt ruộng $z_s(t)$ hạ
với tốc độ $\dot z$, nên một ngưỡng tuyệt đối đo theo một mốc cố định sẽ trôi.

\begin{proposition}[Độ trôi của một ngưỡng tuyệt đối]
\label{prop:subsid}
Trong một chu kỳ cấp tín chỉ dài $T_c$ năm, một ngưỡng tuyệt đối biểu diễn theo một
cao độ gốc cố định sẽ dịch đi $\dot z\,T_c$. Với $\dot z\in[1{,}1;2{,}5]$~cm/năm và
$T_c=5$~năm, độ dịch là $5{,}5$--$12{,}5$~cm, tức $37$--$83\%$ của một ngưỡng
$15$\,cm. Một tiêu chí viết theo tỷ lệ cạn kiệt \eqref{eq:depl} có độ trôi bằng
không, vì $FC$ và $WP$ là tính chất của đất quy về tầng rễ chứ không quy về một cao
độ tuyệt đối.
\end{proposition}

\begin{proof}
Tiêu chí tuyệt đối so một độ sâu đo được với một mốc không đổi; nếu mốc và mặt ruộng
rời nhau với tốc độ $\dot z$ thì ngưỡng hữu hiệu thay đổi $\dot z T_c$ sau $T_c$ năm.
Tiêu chí cạn kiệt so $FC-w_t$ với $(FC-WP)$, cả hai đều được định nghĩa bên trong tầng
rễ và tịnh tiến cứng cùng mặt ruộng, nên tỷ số của chúng bất biến đối với $z_s$.
\end{proof}

Các tốc độ sụt lún là số đo cho đồng bằng này trong~\cite{minderhoud2017}. Hệ quả thực
tiễn là một phương án neo theo tỷ lệ cạn kiệt không cần đo đạc lại các mốc ruộng trong
suốt chu kỳ cấp tín chỉ năm năm, còn phương án dùng độ sâu tuyệt đối thì cần.


% ---------------------------------------------------------------- S5
% S5: ban TIENG VIET cua paper2/blocks/09_cost.tex.
% Chi duoc \input boi paper2_hujos/supplementary_vi.tex.
% \label{sec:cost} GIU LAI vi ban EN van \ref no.
% Chu trong \underbrace dich theo dung ten cot cua Bang S5 (tables_vi/tab_baseline):
% phan cung / ket noi / nhan cong / tham dinh.
\section*{S5. Mô hình chi phí và điểm hòa vốn}
\label{sec:cost}

Giả sử một hợp tác xã phục vụ $A$ héc-ta với $N$ vụ mùa trong đời thiết bị, và chi phí
MRV mỗi héc-ta mỗi vụ phân tích thành
\begin{equation}
C=\underbrace{\frac{n_{\mathrm{dev}}\,p_{\mathrm{dev}}}{A\,N}}_{\text{phần cứng}}
+\underbrace{c_{\mathrm{conn}}}_{\text{kết nối}}
+\underbrace{h\,w}_{\text{nhân công}}
+\underbrace{\frac{c_{\mathrm{audit}}}{A}}_{\text{thẩm định}},
\label{eq:cost}
\end{equation}
trong đó $n_{\mathrm{dev}}$ và $p_{\mathrm{dev}}$ là số lượng và đơn giá thiết bị, $h$
là số giờ công mỗi héc-ta mỗi vụ và $w$ là tiền công nông thôn. Doanh thu mỗi héc-ta là
\begin{equation}
R=\kappa\,\rho ,
\label{eq:revenue}
\end{equation}
với $\kappa$ là số tín chỉ cấp cho mỗi héc-ta và $\rho$ là giá mỗi tín chỉ. Phương án
khả thi về kinh tế khi
\begin{equation}
C\ \le\ \beta R ,\qquad \beta\in(0,1),
\label{eq:breakeven}
\end{equation}
và $\beta=0{,}5$ là mức mà giám sát ngốn hết một nửa doanh thu tín chỉ. Thế
\eqref{eq:revenue} vào ta được lượng tín chỉ tối thiểu để hòa vốn
\begin{equation}
\kappa_{\min}=\frac{C}{\beta\rho}.
\label{eq:kappamin}
\end{equation}

Các tầng khác nhau ở $n_{\mathrm{dev}}$ và, quan trọng hơn, ở chỗ chúng quan sát được
đại lượng nào trong ba đại lượng làm giả được của \eqref{eq:three}. Tầng~0 là tầng vật
lý: $n_{\mathrm{dev}}=0$ ngoài bộ điều khiển, nhật ký được ghi trên điện thoại mà nông
dân đã có sẵn, và các bộ lọc ở Mục~\ref{sec:filters} chạy trên dữ liệu thời tiết công
khai với chi phí biên bằng không. Tầng~1 thêm một đồng hồ lưu lượng tại trạm bơm của hợp
tác xã --- $n_{\mathrm{dev}}=1$ cho toàn bộ diện tích $A$ chứ không phải mỗi thửa một
cái --- và quan sát được tổng lượng nước. Tầng~2 thêm một nút đo mực nước cho mỗi hợp
tác xã, neo lại bản sao số và khôi phục tính hiệu chuẩn mà một hệ MRV dựa trên đo đạc
dù sao cũng cần.

Không tầng nào trong ba tầng đó quan sát được dấu thời điểm, và theo
Mệnh đề~\ref{prop:blind} thì hai tầng đầu bất biến chính xác đối với phép biến đổi sinh
lợi. Vì vậy neo thiết bị ở Mục~\ref{sec:anchor} không phải là một tầng thứ tư tuỳ chọn
mà là điều kiện để bằng chứng được chấp nhận: nó do chính bộ điều khiển đóng máy bơm
sinh ra, nên chỉ thêm việc lưu trữ và truyền đi một bản ghi vốn đã tồn tại, và không
thêm $n_{\mathrm{dev}}$ nào.

Bảng~\ref{tab:baseline} so sánh phương án này với ba kiến trúc mà nó thay thế. Với
$\kappa\in[4,6]$ và $\rho\in[\$15,\$25]$~\cite{rmi2025rice,upstream2026price}, doanh thu
là $R\in[\$60,\$150]$ mỗi héc-ta mỗi vụ, nên $\kappa_{\min}$ tại $\beta=0{,}5$ tương ứng
với một ngân sách giám sát $30$--$75$~USD mỗi héc-ta. Toàn bộ ngăn xếp của bài tốn
$8{,}4$ USD mỗi héc-ta mỗi vụ so với ngân sách $30$--$75$ đó, tức thấp hơn trần từ
$3{,}6$ đến $8{,}9$ lần; vì vậy biên độ được phát biểu thành một dải chứ không phải một
bậc độ lớn. Nó cũng không phải dòng rẻ nhất của Bảng~\ref{tab:baseline}: viễn thám kèm
kiểm tra thực địa tốn $8{,}0$, tức rẻ hơn $0{,}4$. Hai phương án trả lời hai câu hỏi
khác nhau --- phương án rẻ hơn xác nhận rằng một pha khô đã xảy ra nhưng không giữ bản
ghi theo từng lệnh nào để đối chiếu với một nhật ký bị làm giả --- và chi phí tăng thêm
của neo so với một bộ điều khiển đã lắp sẵn là khoảng $1\%$ doanh thu ở mức thận trọng.
Bảng~\ref{tab:baseline} phân tích phép so sánh với ba kiến trúc mà nó thay thế; hạng mục
kết nối được tách riêng vì đó chính là thứ phân biệt một thiết kế mỗi thửa một cảm biến
với một gateway dùng chung.

\begin{remark}
Giá thiết bị trong Bảng~\ref{tab:baseline} là ước tính phục vụ lập kế hoạch, không phải
báo giá. Thứ tự giữa các tầng và bất đẳng thức hòa vốn \eqref{eq:breakeven} không phụ
thuộc vào chúng; chỉ có biên độ tuyệt đối là phụ thuộc. Cần thay các ước tính bằng báo
giá trước bất kỳ quyết định mua sắm nào.
\end{remark}


\begin{table}[!h]
\centering
\caption{Chi phí MRV mỗi héc-ta mỗi vụ cho bốn kiến trúc. Giám sát chiếm phần lớn chi phí
MRV, nên kiến trúc tránh được cảm biến chuyên dụng thắng ở phần nhân công chứ không phải
ở phần cứng.}
\label{tab:baseline}
\footnotesize
% Sinh TU DONG tu tables/tab_baseline.tex (ban tieng Anh, von sinh tu outputs/*.csv)
% boi experiments/make_tables_hujos_vi.py. KHONG sua tay.
% Nhan dich qua bang tra cuong che; so chi doi dinh dang sang kieu Viet Nam (dau phay thap phan).
% Da assert: cung so dong, cung so vach, cung tap chu so voi ban EN.

% Sinh TU DONG tu outputs/eval_theta_*.csv boi experiments/make_tables_anchored.py
% KHONG sua tay: moi con so deu truy nguon duoc ve CSV.

\fittocol{%
\begin{tabular}{@{}>{\raggedright\arraybackslash}p{3.00cm}>{\raggedright\arraybackslash}p{1.10cm}>{\raggedright\arraybackslash}p{1.30cm}>{\raggedright\arraybackslash}p{1.00cm}>{\raggedright\arraybackslash}p{0.85cm}>{\raggedright\arraybackslash}p{1.10cm}>{\raggedright\arraybackslash}p{4.00cm}@{}}
\toprule
Kiến trúc & Phần cứng & Kết nối & Nhân công & Thẩm định & Tổng \$/ha & Ghi chú \\
\midrule
Ống quan trắc thủ công kèm ảnh chụp & 0,5 & 0,0 & 8,0 & 4,0 & \textbf{12,5} & thực hành hiện hành theo Quyết định 4801 \\
Cảm biến mực nước mỗi thửa & 11,3 & 2,0 & 1,5 & 4,0 & \textbf{18,8} & chính xác, khó mở rộng cho nông hộ nhỏ \\
Viễn thám kèm kiểm tra thực địa & 0,0 & 0,0 & 2,0 & 6,0 & \textbf{8,0} & chu kỳ chụp 2 đến 4 ngày \\
Nhật ký được chứng thực, bài này & 3,6 & 1,5 & 0,3 & 3,0 & \textbf{8,4} & dấu thời gian của gateway là thứ làm nó được chấp nhận \\
\bottomrule
\end{tabular}
}

\end{table}

% ---------------------------------------------------------------- S6
\section*{S6. Cấu hình thực nghiệm}

\begin{table}[!h]
\centering
\caption{Cấu hình thực nghiệm. Mọi giá trị đều tái lập được từ mã nguồn đã công bố và dữ
liệu ERA5 công khai.}
\label{tab:setup}
\footnotesize
% Sinh TU DONG tu tables/tab_setup.tex (ban tieng Anh, von sinh tu outputs/*.csv)
% boi experiments/make_tables_hujos_vi.py. KHONG sua tay.
% Nhan dich qua bang tra cuong che; so chi doi dinh dang sang kieu Viet Nam (dau phay thap phan).
% Da assert: cung so dong, cung so vach, cung tap chu so voi ban EN.

% Sinh TU DONG tu outputs/eval_theta_*.csv boi experiments/make_tables_anchored.py
% KHONG sua tay: moi con so deu truy nguon duoc ve CSV.

\fittocol{%
\begin{tabular}{@{}>{\raggedright\arraybackslash}p{2.85cm}>{\raggedright\arraybackslash}p{5.60cm}@{}}
\toprule
Thành phần & Giá trị \\
\midrule
Thời tiết & ERA5 theo giờ qua kho lưu trữ Open-Meteo, không có giá trị thiếu \\
Trạm & 3: Cà Mau, Cần Thơ, Sóc Trăng \\
Năm khí hậu & 2024, 2025 (2024 khan nước, 2025 thừa nước) \\
Cửa sổ & từ ngày 1 tháng Ba đến ngày 29 tháng Năm (2\,160 h), mùa khô đồng bằng \\
Loại đất được lấy mẫu & 5 loại trải theo FAO-56 Bảng 2 \\
Trục khác biệt thực tế & 7: đất, giai đoạn $K_c$, $\eta$, mưa tại ruộng, thấm bờ, hành vi ghi chép, $w_0$ \\
Hiệu suất tưới $\eta$ & 0,60--0,85 (nước tới rễ / nước bơm) \\
Chế độ $K_c$ & cố định 1,05 (như đã công bố) và theo giai đoạn 1,05/1,20/0,90 \\
Nhiễu ghi chép & quên 5\%, làm tròn đến 5\,mm, trễ một ngày 35\%, sai độ sâu 15\% trong $\pm$35\% \\
Biến thể gian lận & 7 (ba ngây thơ, bốn đối kháng) \\
Tập bất định $|\Theta|$ & 30 = 5 loại đất $\times$ 2 chế độ $K_c$ $\times$ 3 hiệu suất \\
Trạng thái đầu $|w_0|$ & 5, trải cạn kiệt từ 0 đến 0,90 \\
Số lần phát lại mỗi nhật ký & 24 ($\Theta$ rút gọn) và 150 ($\Theta$ đầy đủ) \\
Số nhật ký được kiểm & 281 \\
\bottomrule
\end{tabular}
}

\end{table}

% ---------------------------------------------------------------- S7
\section*{S7. Hai chi tiết triển khai của phép giao bền vững}

Có hai chi tiết triển khai của $V_{\mathrm{rob}}$ là quan trọng, và làm sai chi tiết nào
cũng âm thầm phá vỡ Mệnh đề~\ref{prop:noalarm} của bài chính. Một thế giới không sinh ra
vi phạm nào phải đóng góp \emph{tập rỗng} vào phép giao; bỏ qua các thế giới không kích
hoạt là tính $\bigcap$ chỉ trên các thế giới có kích hoạt, tức một lượng từ tồn tại mặc
áo lượng từ phổ dụng. Và $\mathcal{W}$ phải trải hết dải mà một khai báo còn có thể hợp
lệ: với $\mathcal{W}=\{0;0{,}25\}$ thì độ sâu chấp nhận được lớn nhất ở bất kỳ đâu trong
$\Theta$ là $22{,}5$\,mm, nên mọi khai báo vượt mức đó đều ``có tội trong mọi thế giới''
và phép giao trở nên rỗng nghĩa --- nó kích hoạt trên $102$ trong $150$ thế giới trong khi
vẫn tồn tại $48$ thế giới vô tội. Mở rộng $\mathcal{W}$ thành
$\{0;0{,}25;0{,}5;0{,}75;0{,}9\}$ khôi phục đúng ngữ nghĩa dự kiến.

% ---------------------------------------------------------------- S8
\section*{S8. Quan hệ với các công trình giám sát trước đây (toàn văn)}

Giám sát AWD đã được tiếp cận từ ba hướng. \emph{Gắn thiết bị}: một cảm biến mực nước
không dây giá rẻ cho lúa đồng bằng theo AWD~\cite{delacruz2022}, giải quyết được bài toán
đo lường nhưng nhân chi phí thiết bị theo số thửa. \emph{Viễn thám}: nhận diện việc áp
dụng AWD ở đồng bằng sông Cửu Long~\cite{lovell2019}, giám sát trạng thái nước bằng phân
tách công suất tán xạ ALOS-2~\cite{islam2026alose}, và lập bản đồ học máy ở quy mô vùng
cho các thực hành tiết kiệm nước~\cite{mlrice2025}; các hướng này phủ diện rộng với chi
phí thấp nhưng hai đến bốn ngày mới chụp lại một lần nên không phân giải được pha khô
$72$\,h vốn quyết định điều kiện cấp tín chỉ. \emph{Cấp chương trình}: chương trình một
triệu héc-ta lúa chất lượng cao, phát thải thấp~\cite{qd1490}, đến năm 2026 đã đưa hơn
$350\,000$ héc-ta vào canh tác~\cite{mae2026progress}, và MRV được xác định là mắt xích
còn thiếu của lúa phát thải thấp ở Đông Nam Á~\cite{irri2026mrv}.

Về phía nông học, phân tích tổng hợp định lượng tác động của AWD lên nước và năng
suất~\cite{carrijo2017}, còn kinh tế học về việc chấp nhận~\cite{lampayan2015} cùng các
rào cản kỹ thuật và thể chế khi mở rộng~\cite{enriquez2021} giải thích vì sao giám sát
gắn thiết bị chưa lan rộng. Không hướng nào trong số đó đánh giá \emph{độ tin cậy của
chính bản ghi hành động}, vốn là câu hỏi đặt ra ở đây.

Gateway cung cấp dấu thời điểm được neo theo thiết bị chính là bộ điều khiển của công
trình đồng hành~\cite{twingate2026}, trong đó việc phát lệnh chạy trên một kênh nhiễu
chùm loại Gilbert--Elliott~\cite{gilbert1960,elliott1963}. Các kết quả lập lịch trước đây
của chính chúng tôi cho mạng IoT nông nghiệp an toàn tới hạn cũng xử lý cùng chế độ
truyền đó~\cite{rabs,cawvou}, nhưng không công trình nào thiết lập một bản ghi bằng
chứng, vốn là chức năng được dựa vào ở đây. Vì bộ điều khiển đó đã ký mọi lệnh nó phát
ra, neo chỉ thêm việc lưu trữ và truyền đi một bản ghi vốn đã tồn tại --- đó chính là
điều giữ cho chi phí biên của việc chứng thực gần bằng không.

% ---------------------------------------------------------------- S9
\section*{S9. Bàn luận mở rộng}

\subsection*{S9.1 Vì sao công bố ngưỡng không phải là vấn đề}

Có thể kết luận rằng nên giữ bí mật $\varepsilon$. Không nên, vì hai lý do. Một kiểm toán
viên phải tái lập được một phán quyết, và điều đó đòi hỏi tham số được công bố. Và giữ bí
mật cũng không ích gì: theo Mệnh đề~\ref{prop:blind} của bài chính, tấn công dịch thời
điểm giữ nguyên tổng lượng và từng độ sâu riêng lẻ, nên nó không cần biết $\varepsilon$ để
nằm dưới giới hạn thấm --- nó chỉ cần tránh dịch chuyển nước, điều mà bất kỳ nông dân nào
cũng nhìn thấy. Lỗ hổng có tính cấu trúc, không phải hệ quả của việc công bố.

\subsection*{S9.2 Quan hệ với viễn thám (toàn văn)}

Radar vệ tinh là phương án giá rẻ hiển nhiên, nên thay vì tranh luận về nó chúng tôi đo
nó. Bảng~\ref{tab:benchmark} của bài chính chấm điểm mọi phương án bằng một đại lượng:
trong số $150$ pha khô đủ điều kiện cấp tín chỉ có thật trên $60$ ruộng mô phỏng, phương án
đó xác nhận được bao nhiêu. Một pha chỉ được tính là xác nhận được nếu có ít nhất một lần
quan sát rơi vào trong nó, vì một pha $72$\,h nằm trọn giữa hai lần chụp thì không thể có
bằng chứng nào --- điều đó làm nông dân mất một tín chỉ ngay cả khi thực hành được thực
hiện đúng. Sentinel-1 một vệ tinh, chu kỳ mười hai ngày~\cite{clauss2017s1}, xác nhận
$90/150$ ($60\%$); cả chòm A+B với chu kỳ sáu ngày~\cite{hashemi2022s1} đạt $124/150$
($83\%$), và vẫn bỏ sót một pha trong mỗi sáu pha dù chi phí thiết bị bằng không. Một ống
đo mực nước liên tục xác nhận cả $150$, với mỗi thửa một thiết bị~\cite{delacruz2022}. Bản
ghi lệnh có chữ ký cũng xác nhận cả $150$ --- vì các pha được định nghĩa bằng những khoảng
giữa hai lần đóng mở, mà bộ điều khiển vốn đã đóng dấu thời điểm --- và không cần thêm
thiết bị nào, vì chính bộ điều khiển là thứ đã mở máy bơm.

Có hai điểm cần thận trọng kèm theo phép so sánh đó. Cột xác nhận đo khả năng \emph{kiếm}
được tín chỉ, chứ không đo khả năng phát hiện gian lận: ống đo và bản ghi lệnh hòa nhau ở
$100\%$ tại đây nhưng khác nhau hoàn toàn về chi phí và về những gì chúng chứng minh được
đối với một nhật ký đã bị sửa. Còn phương án chụp ảnh cộng tự khai báo đạt $0\%$ không phải
vì nó vô giá trị mà vì nó không sinh ra chuỗi quan sát nào để đối chiếu; mọi xác nhận theo
nó đều phải đến từ một chuyến đi thực địa.

Vì vậy radar giữ đúng vai trò mà VM0051 Phụ lục 4 giao cho nó: một phép kiểm chéo theo
từng thửa làm giảm tần suất xác minh thực địa~\cite{vm0051}.

% ---------------------------------------------------------------- S10
\section*{S10. Các lỗi của bộ khung thử nghiệm và chi tiết thiết kế}

\subsection*{S10.1 Hành vi ghi chép theo từng sự kiện chứ không theo từng giờ}

Việc làm tròn và trễ được áp theo từng \emph{sự kiện}, không theo từng giờ: áp mức làm tròn
$5$\,mm cho một sự kiện $40$\,mm rải trên tám giờ sẽ làm tròn mỗi mảnh $5$\,mm về không và
âm thầm xóa mất bản ghi, một ảo ảnh đã từng sinh ra các nhật ký rỗng trong phiên bản đầu
của bộ khung này.

\subsection*{S10.2 Hai lỗi bộ khung phải sửa}

Phải sửa hai lỗi của bộ khung trước khi các con số được báo cáo có ý nghĩa gì, và cả hai
đều được ghi lại trong mã nguồn vì chúng rất dễ bị đưa trở lại. Thứ nhất, một trình sinh
gian lận không tìm được nước đi cải thiện sẽ trả lại nhật ký trung thực \emph{ban đầu}, nên
nếu đếm vô điều kiện thì số ``tấn công bị bắt'' âm thầm gộp cả những ca không hề có tấn
công; vì vậy mọi dòng đối kháng đều mang một cờ \texttt{changed} và tất cả thống kê tấn
công đều được tính có điều kiện trên cờ đó. Thứ hai, các phép so sánh phải dùng mốc của
chính biến thể đó: nhật ký trung thực điều kiện giấy được sinh ra bởi một lần gọi riêng,
và việc chấm nó theo bản ghi điều kiện ruộng của một lần hiện thực hóa đất khác đã ``bắt''
sai $56$ trong $60$ nhật ký trung thực.

\subsection*{S10.3 Tổng lượng là chỉ báo chẩn đoán, không phải mục tiêu}

Một cấu hình tối đa hóa tổng lượng khai báo không tự động sinh lợi: một đối thủ phiên bản
sớm thổi phồng khai báo thêm $344$\,mm đã kéo $n_{\mathrm{dry}}$ về không và sẽ chẳng kiếm
được gì, nên tổng lượng chỉ được báo cáo như một chỉ báo chẩn đoán.

\subsection*{S10.4 Ảo ảnh tự nhất quán}

Tầng được kiểm toán từng sinh ra các quỹ đạo ``thật'' bằng \texttt{SoilProfile} mặc định,
dữ liệu ERA5 không nhiễu, tưới là một xung một giờ đúng bằng $\{20,40\}$\,mm, ngưỡng đúng
bằng $0{,}80$, và một nông dân luôn ghi đúng độ sâu vào đúng giờ. Vì vậy cả $0/60$ ca kết
tội oan lẫn $60/60$ ca phát hiện đều là phép đo tính tự nhất quán: bộ phát hiện đang chấm
chính trình sinh đã tạo ra các quỹ đạo ``thật'' mà nó kiểm toán. Trong thiết kế hiện tại,
thửa ruộng được sinh bởi một tiến trình riêng biệt, khác bản sao số theo bảy trục liệt kê
ở Bảng~S5.

% ---------------------------------------------------------------- S11
\section*{S11. Chứng minh Mệnh đề 1 và Mệnh đề 2}

\subsection*{Chứng minh Mệnh đề 1 (sự mù chính xác của các tầng phòng thủ tổng hợp)}

$U(L)=\sum_i u_i$ là một hàm đối xứng của đa tập hợp, nên nó bất biến đối với mọi phép gán
lại độ sâu cho các giờ; $M$ là tính chất của thửa ruộng chứ không phải của nhật ký; do đó
$G=M-U$ không đổi. Tương tự, giới hạn thấm (\eqref{eq:v5} của bài chính) được đánh giá
theo từng sự kiện và chỉ phụ thuộc $(u_i,\tau_i)$, không bao giờ phụ thuộc $t_i$ hay
$\hat w$, nên số lần vi phạm là một hàm chỉ của đa tập hợp.

\subsection*{Chứng minh Mệnh đề 2 (lợi ích tín chỉ và sự lọt lưới là cùng một hành động)}

Phần (i) là theo định nghĩa: tập $\mathcal{A}$ không đòi hỏi khai báo nào bên trong khoảng,
nên bỏ đi một khai báo sẽ kéo dài hoặc tạo ra một pha đạt điều kiện. Với phần (ii), hoãn
một khai báo đến $t+\delta$ nghĩa là thùng chứa tiếp tục mất nước thêm $\delta$ giờ mà không
có nước vào, nên $\hat w_{t+\delta}^{-}\le\hat w_t^{-}$ và $c_{t+\delta}\ge c_t$; cận trên
của độ tăng phần thiếu hụt suy ra từ chỗ lượng mất mỗi giờ ít nhất bằng trung bình
$\overline{\Delta}$ cho tới khi đạt sức chứa ẩm đồng ruộng, sau đó $c$ bão hòa tại
$FC-\hat w$. Vì bộ lọc vượt sức chứa chỉ kích hoạt khi $u_i>c_{t_i}+\varepsilon$ và $u_i$
không đổi, một $c_{t_i}$ lớn hơn chỉ có thể làm giảm tập vi phạm.

% ---------------------------------------------------------------- S12
\section*{S12. Chứng minh Mệnh đề 4 (tính đầy đủ của neo)}

Vế thứ nhất chính là định nghĩa của phán quyết neo: bất kỳ sự bất đồng nào giữa các tập giờ,
hoặc giữa độ sâu tại một giờ chung, đều kích hoạt phán quyết giả mạo. Với vế thứ hai, tổng
lượng, từng độ sâu riêng lẻ và các dấu thời điểm đều là hàm của tập giờ và độ sâu theo giờ,
nên sự trùng khớp trên hai thứ đó cố định cả ba, và không có sự giả mạo nào đối với ba đại
lượng đã xảy ra.

% ---------------------------------------------------------------- S13
\section*{S13. Bộ lọc điểm héo: toàn bộ dẫn xuất}

Tầng này từng gồm một bộ lọc vật lý thứ ba, đánh dấu việc cạn kiệt xuống dưới điểm héo là
bất khả thi về mặt vật lý. Nó vô hiệu theo hai nghĩa. Ngưỡng của nó được đặt ở $1{,}02$
trong khi chỉ báo được báo cáo $\hat d_t=\min(1,\max(0,\cdot))$ bị cắt về $[0,1]$, nên điều
kiện không thể thỏa mãn ngay từ cách dựng: $0$ lần kích hoạt trong $281$ nhật ký. Chúng tôi
sửa nó bằng cách đánh giá độ cạn kiệt chưa cắt ngọn, rồi đo xem bộ lọc đã sửa có kích hoạt
được hay không. Trên ba trạm, hai năm khí hậu, năm loại đất và năm trạng thái đầu --- tức
$150$ cấu hình --- với nhật ký rỗng, nghĩa là điều kiện thời tiết khô hạn nhất có thể, độ
cạn kiệt chưa cắt ngọn đạt đỉnh $0{,}9987$ (Bảng~S3). Một ruộng lúa vụ khô ở đồng bằng
không chạm tới điểm héo. Vì vậy bộ lọc này là một thiết bị \emph{an toàn} bảo đảm không một
trạng thái vô lý nào được chấp nhận; nó không phải bằng chứng về năng lực phát hiện, và tầng
này phải được mô tả là có hai bộ lọc vật lý hoạt động và một bộ lọc trơ.

% ---------------------------------------------------------------- S14
\section*{S14. Hai hiệu chỉnh đối với biên độ sâu chấp nhận được}

Thứ nhất, \eqref{eq:umax} của bài chính dùng $\max$ chứ không dùng một hiệu số: một thửa
ruộng khô ban đầu ($w_0=0{,}90$) chấp nhận những khai báo \emph{lớn hơn} một thửa ruộng ướt,
còn biểu thức ban đầu $(d^{*}-w_0)\,\mathrm{span}+\varepsilon$ cho ra $43{,}5$\,mm với
$\mathrm{span}=70$\,mm và $w_0=0{,}25$, một giá trị không tương ứng với trạng thái vật lý
nào, vì phần thiếu hụt hữu hiệu tại giờ đó hoặc là $0{,}25\,\mathrm{span}$ (lúc đầu) hoặc là
$d^{*}\mathrm{span}$ (khi ruộng đã khô). Thứ hai, biên này áp theo \emph{từng sự kiện} chứ
không theo từng vụ: thùng chứa rút liên tục qua thấm sâu và bốc thoát hơi nước, nên tổng
khai báo cả vụ không bị chặn bởi $u_{\max}$ và một đối thủ phát ra nhiều khai báo dưới
ngưỡng có thể thổi phồng tổng lượng mà không có giới hạn.

% ---------------------------------------------------------------- S15
\section*{S15. Vì sao mốc ghép cặp mới là mốc trung thực}

Các biến thể giả mạo chỉ được sinh ở những quỹ đạo mang đủ số lệnh để thao tác, tức từ ba
lệnh trở lên, và những quỹ đạo đó vốn đã khô hơn trung bình: mốc trung thực của chúng mang
$4{,}87$ pha khô cấp được tín chỉ so với $3{,}23$ của toàn bộ nhật ký trung thực. Vì vậy một
phép so sánh không ghép cặp với mốc rộng hơn sẽ quy phần khô hạn vốn có cho cuộc tấn công và
thổi phồng mức tăng lên xấp xỉ gấp đôi. Mọi con số trong Bảng~\ref{tab:credit} của bài chính
đều được tính trên tập con ghép cặp, cùng trạm, cùng năm, cùng hạt giống và cùng loại đất,
có điều kiện trên cờ \texttt{changed}.

% ---------------------------------------------------------------- S16
\section*{S16. Chín loại nhật ký}

Ba loại là giả mạo ngây thơ: thêm sự kiện vào ngày mưa, thổi phồng một phần ba số mục lên
bằng toàn bộ biên độ chứa, và xóa một cửa sổ khai báo dài $96$\,h. Bốn loại là đối kháng và,
khác với các loại ngây thơ, được dựng bằng một bộ tối ưu chứ không theo một công thức cố
định: gộp khai báo về lưới $24$\,h và lưới $96$\,h; một đối thủ duy lý chọn cửa sổ gộp nào
làm tối đa $n_{\mathrm{dry}}$ với điều kiện giữ nguyên tổng lượng và tôn trọng giới hạn
thấm; và phép leo đồi theo dịch thời điểm, tìm các độ lệch nguyên trong $[-36,36]$ và chỉ
chấp nhận một nước đi nếu tổng lượng được giữ nguyên từng bít, mọi độ sâu đều nằm dưới giới
hạn thấm, và $n_{\mathrm{dry}}$ tăng thực sự. Hai loại trung thực còn lại, điều kiện giấy và
điều kiện ruộng, là các mốc so sánh.

% ---------------------------------------------------------------- S17
\section*{S17. Cái giá của tính bền vững và sự nhất quán giữa các tập bất định}

\subsection*{S17.1 Mức mất mát năng lực phát hiện đo được}

Thay mô hình điểm bằng phép giao trên $\Theta$ không miễn phí, và mức mất mát phải được báo
cáo. Trên chín nhật ký gian lận bỏ sót thật sự bị sửa, mô hình điểm bắt được $9/9$ trong khi
quy tắc bền vững bắt được $4/9$ ($44\%$); với $\Theta$ đầy đủ các con số là $4/4$ và $2/4$
($50\%$). Cùng một sự đánh đổi chạy theo cùng hướng với ngưỡng tỉ lệ, và vì cùng một lý do:
một tầng phòng thủ từ chối kết tội chỉ dựa trên sức mạnh của một cấu hình mô hình thì cũng từ
chối một số ca dương tính thật. Giả mạo ngây thơ vẫn ở mức $100\%$ theo cả hai. Chúng tôi
coi đây là phía đúng của sự đánh đổi, vì một ca kết tội oan rơi xuống một nông dân không có
cách nào kháng cáo còn một ca giả mạo bị bỏ lọt thì rơi xuống một cơ quan xác minh vốn vẫn
lấy mẫu, nhưng đó là sự mất mát năng lực phát hiện chứ không phải một cải thiện Pareto.

\subsection*{S17.2 Tập rút gọn và tập đầy đủ cho cùng kết quả}

Mọi kết luận đều đúng với cả $\Theta$ rút gọn, gồm $24$ lần phát lại mỗi nhật ký, $281$ nhật
ký và $10$ hạt giống ngẫu nhiên, lẫn $\Theta$ đầy đủ, gồm $150$ lần phát lại, $85$ nhật ký và
$3$ hạt giống ngẫu nhiên: kết tội oan $0/60$ và $0/18$, tính đầy đủ của neo $116/116$ và
$37/37$, lọt lưới dịch thời điểm ở tầng bền vững $8/9$ và $3/3$, thổi phồng tín chỉ ghép cặp
$+33\%$ ($t=5{,}29$) và $+38\%$ ($t=3{,}46$). Tập rút gọn được dùng cho các lần chạy lớn vì
nó chứa các cực trị của mọi trục, tức nơi một thế giới vô tội dễ tồn tại nhất; tập đầy đủ xác
nhận rằng không một cấu hình nằm bên trong nào làm thay đổi phán quyết.

% ---------------------------------------------------------------- S18
\section*{S18. Vì sao neo không đưa việc kết tội oan trở lại (toàn văn)}

Mệnh đề~\ref{prop:complete} của bài chính mạnh tới mức thoạt trông có vẻ nguy hiểm: mọi sai
lệch đều bị coi là giả mạo, kể cả việc quên ghi và làm tròn mà bài chính đã chỉ ra là rất phổ
biến. Cách hóa giải nằm ở chỗ neo không đối chiếu một nhật ký với một \emph{mô hình}; nó đối
chiếu một bản chép với \emph{nguồn}. Một nông dân trung thực quên một mục không bị kết tội
gian lận, vì hệ thống lưu trữ hồ sơ là gateway chứ không phải trí nhớ của nông dân: khai báo
đó đơn giản được bổ sung từ $\mathcal{G}$. Về mặt vận hành, neo biến nhật ký từ \emph{bằng
chứng mà nông dân phải ghi cho đúng} thành \emph{một góc nhìn của bản ghi mà thiết bị đã
giữ sẵn}. Khi đó việc kết tội oan đòi hỏi gateway phải bất đồng với chính nó, vốn là một câu
hỏi về toàn vẹn thiết bị chứ không phải về hành vi, và đó là lý do bảo đảm này giữ ở mức
$0/60$ trong khi các tầng dựa trên hành vi giữ ở mức $26{,}7\%$ và $5{,}0\%$.

% ---------------------------------------------------------------- S19
\section*{S19. Cuộc tấn công lấy đi thứ gì}

Cuộc tấn công không lấy đi nước. Tổng lượng nước bơm không đổi, nên nông dân trả vẫn chừng
đó và khí quyển vẫn thấy chừng đó mê-tan. Thứ bị lấy đi là \emph{lịch sử}: tín chỉ được cấp
cho một mẫu pha khô đã không xảy ra. Theo VM0051, đơn vị cấp tín chỉ chính xác là mẫu đó, nên
thiệt hại là có thật dù không có nguồn lực vật lý nào dịch chuyển.

% ---------------------------------------------------------------- S20
\section*{S20. Thiết kế thực nghiệm: toàn bộ chi tiết}

\subsection*{S20.1 Gỡ bỏ tính vòng tròn, theo từng trục}

Thửa ruộng được sinh bởi một tiến trình riêng biệt, khác bản sao số theo bảy trục, mỗi trục
đều có một nguồn đã công bố (Bảng~S5): thành phần cơ giới của đất lấy từ năm cấp đất của
FAO-56 Bảng 2 có mặt ở đồng bằng; $K_c$ phân theo giai đoạn đúng Bảng 12 chứ không cố định ở
$1{,}05$; hiệu suất tưới $\eta\in[0{,}60;0{,}85]$, nên nông dân khai lượng nước \emph{đã bơm}
trong khi chỉ $\eta u$ tới được tầng rễ; mưa tại ruộng lấy bằng $0{,}7$--$1{,}4$ lần ô lưới
cộng nhiễu Gauss; thấm bờ $0$--$0{,}02$\,mm\,h$^{-1}$; hành vi ghi chép với $5\%$ quên ghi,
làm tròn đến $5$\,mm, $35\%$ khả năng trễ một ngày và $15\%$ sai độ sâu trong phạm vi
$\pm35\%$; và độ cạn kiệt ban đầu $w_0$ lấy từ $[0;0{,}45]$ thay vì được giả định là sức chứa
ẩm đồng ruộng. Việc làm tròn và trễ được áp theo từng \emph{sự kiện} chứ không theo từng giờ
(S10.1).

\subsection*{S20.2 Phần còn lại được của tuyên bố ban đầu}

Tiền đề thúc đẩy công trình này vốn hấp dẫn: một lệnh tưới tự thân đã là một sự kiện số, nên
bằng chứng MRV lẽ ra gần như miễn phí. Tiền đề đó vẫn đứng vững, nhưng không ở dạng được
tuyên bố ban đầu. Nó đứng vững vì bộ điều khiển phát ra lệnh cũng có thể ký lệnh đó, và chính
bản ghi có chữ ký là thứ làm cho bằng chứng được chấp nhận --- chứ không phải vì một mô hình
thủy văn có thể phát hiện một lời khai dối. Ở chỗ bản thảo ban đầu tuyên bố rằng vật lý chứng
thực nhật ký, phát biểu đúng là vật lý \emph{ràng buộc} nhật ký: nó loại trừ những độ sâu mà
không loại đất nào chứa được và những trạng thái mà không thửa ruộng nào đạt tới, điều đó
đáng có, và nó im lặng về đúng phép biến đổi duy nhất sinh lợi. Giữ bí mật $\varepsilon$ cũng
không bịt được lỗ hổng: cuộc tấn công không cần biết ngưỡng, và dù sao một kiểm toán viên
cũng phải tái lập được phán quyết (S9.1).

% ---------------------------------------------------------------- S21
\section*{S21. Các giới hạn (toàn văn)}

\begin{enumerate}
\item \emph{Nông dân mô phỏng.} Các hành vi giả mạo được tiêm vào những quỹ đạo tham chiếu mô
phỏng. Cần một thử nghiệm thực địa để đo phản ứng hành vi, cụ thể là liệu nông dân có thay
đổi việc bơm \emph{thực tế} của họ một khi biết rằng neo tồn tại hay không --- một kiểu tấn
công mà mô hình của chúng tôi không biểu diễn được, vì nó thay đổi $M$ chứ không thay đổi $L$.
\item \emph{Neo giả định một gateway trung thực.} Mệnh đề~\ref{prop:complete} của bài chính là
tính đầy đủ đối với việc giả mạo nhật ký, không phải đối với một bộ điều khiển đã bị chiếm
quyền. Toàn vẹn thiết bị, việc giữ khóa và chống phát lại được giả định từ công trình đồng
hành và không được thiết lập lại ở đây.
\item \emph{Không áp dụng cho tưới thủ công.} Nơi một cửa cống được mở bằng tay thì không có
bản ghi lệnh, nên neo không khả dụng. Điều này loại trừ một phần đáng kể số nông hộ ở đồng
bằng và giới hạn quần thể mà phương án có thể phục vụ.
\item \emph{Cuộc tấn công chỉ được minh họa trong năm khan nước.} Trong năm thừa nước (2025),
nông dân cần tưới ít đến mức phần lớn nhật ký mô phỏng rơi xuống dưới mức tối thiểu ba lệnh,
mức mà một tấn công thời điểm mới có thể dựng được: trong $37$ nhật ký thật sự bị sửa ở
$\Theta$ đầy đủ thì cả $37$ đều thuộc năm 2024. Chính vách ngăn giữa hai năm đó đã là một
phát hiện --- rủi ro gian lận MRV tập trung ở nơi khan nước --- nhưng nó cũng có nghĩa là các
kết quả tấn công chưa phải bằng chứng về một mùa mưa. Ở $\Theta$ rút gọn có mười hai nhật ký
bị sửa thuộc năm 2025 và chúng hành xử giống nhau, điều đó gợi ý chứ chưa kết luận được.
\item \emph{$\Theta$ là một phán đoán chuyên môn.} Năm loại đất, hai chế độ $K_c$, ba mức hiệu
suất và năm trạng thái đầu bao được đồng bằng nhưng không vét hết nó. Một thửa ruộng nằm ngoài
$\Theta$ sẽ mất bảo đảm của Mệnh đề~\ref{prop:noalarm} trong bài chính, và đó chính xác là
tình huống mà ngưỡng tỉ lệ cùng phán quyết chờ xác minh tồn tại để hấp thụ.
\item \emph{Giá thiết bị là ước tính phục vụ lập kế hoạch, không phải báo giá.} Thứ tự giữa các
tầng và bất đẳng thức hòa vốn không bị ảnh hưởng; biên độ tuyệt đối thì có.
\item \emph{Không định lượng phát thải.} Tầng này sinh ra dữ liệu hoạt động. Việc ghép nó với
một mô hình mê-tan và các hệ số phát thải được hiệu chuẩn tại chỗ là một công trình riêng,
cũng như giả định của bảng chi phí rằng tín chỉ tăng theo $n_{\mathrm{dry}}$.
\end{enumerate}

\bibliographystyle{vancouver}
\bibliography{refs}

\end{document}
